% ============================================================================
% MÓDULO 7: MACHINE LEARNING PARA BUSINESS INTELLIGENCE
% ============================================================================
\chapter{Machine learning para Business Intelligence}
\label{cap:ml}

\begin{objetivos}
\begin{itemize}
  \item Explicar con palabras de negocio qué es el \emph{machine learning}, qué
        problemas resuelve y cuáles no.
  \item Distinguir aprendizaje supervisado (regresión y clasificación) y no
        supervisado (clustering y reglas de asociación), y elegir el algoritmo
        adecuado para cada pregunta de negocio.
  \item Aplicar el flujo de trabajo completo: separar datos de entrenamiento y
        prueba, entrenar, evaluar con las métricas correctas y evitar el
        sobreajuste y la fuga de datos.
  \item Segmentar clientes con RFM y \emph{k-means}, y darle un nombre de
        negocio a cada segmento.
  \item Predecir ventas con regresión lineal y \emph{random forest}, y medir el
        error con MAE, RMSE, MAPE y R².
  \item Predecir qué clientes van a abandonar (\emph{churn}) con regresión
        logística, árboles de decisión y \emph{random forest}; leer una matriz
        de confusión, una curva ROC y elegir el umbral según el costo.
  \item Pronosticar la demanda mensual con modelos de series de tiempo
        (ingenuo, ETS y ARIMA) y comunicar la incertidumbre.
  \item Descubrir qué productos se compran juntos (análisis de canasta) para
        campañas de venta cruzada.
  \item Guardar un modelo, usarlo para puntuar clientes nuevos, vigilar que siga
        funcionando y usarlo de forma ética.
\end{itemize}
\end{objetivos}

Es lunes por la mañana y la directora comercial te escribe tres preguntas:
\emph{``¿Qué clientes están a punto de irse y a quiénes llamamos primero?'',
``¿Cuánto vamos a vender en diciembre para pedir el inventario a tiempo?'' y
``¿A qué clientes les mandamos la campaña VIP?''}. Hasta ahora, con lo que
aprendiste en los módulos anteriores, sabes responder muy bien \emph{qué pasó}:
cuánto se vendió, en qué tienda, qué producto fue el más rentable. Estas tres
preguntas son distintas: piden \emph{anticipar} lo que va a pasar y
\emph{decidir} qué hacer. Ahí entra el \termino{machine learning} (aprendizaje
automático).

En este módulo vas a construir, paso a paso y con código comentado, los modelos
que más se usan en Business Intelligence: segmentación de clientes, predicción
de ventas, predicción de abandono, pronóstico de demanda y análisis de canasta.
No necesitas matemáticas avanzadas: te explicaremos cada modelo con analogías
de negocio, y nos concentraremos en lo que de verdad importa en una empresa:
\textbf{hacer la pregunta correcta, evaluar honestamente el modelo y convertir
el resultado en una decisión}.

\begin{concepto}{Las cuatro preguntas de la analítica}
\begin{itemize}
  \item \textbf{Descriptiva:} ¿qué pasó? (reportes, tableros; Módulos~\ref{cap:manipulacion}
        y~\ref{cap:visualizacion}).
  \item \textbf{Diagnóstica:} ¿por qué pasó? (análisis exploratorio, Módulo~\ref{cap:eda}).
  \item \textbf{Predictiva:} ¿qué va a pasar? (machine learning, este módulo).
  \item \textbf{Prescriptiva:} ¿qué debemos hacer? (combinar la predicción con
        costos y reglas de negocio; también lo verás aquí, por ejemplo al
        elegir a qué clientes llamar).
\end{itemize}
\end{concepto}

Una advertencia honesta desde el principio: los datos del curso (por
ejemplo, \archivo{ventas\_retail.csv} o \archivo{transacciones.csv}) se generaron
\emph{al azar}. Sirven perfectamente para la segmentación de clientes, porque
ahí solo queremos \emph{describir} grupos, pero no tienen patrones que un modelo
predictivo pueda aprender. Por eso, para predicción de ventas, abandono de
clientes, series de tiempo y canastas vamos a \textbf{simular datos con
relaciones conocidas} (como si fueran los datos de un cliente real). Tiene una
ventaja didáctica enorme: como sabemos la ``verdad'' con la que se generaron,
podremos comprobar si el modelo la descubre. Y en la
sección~\ref{sec:m7-sin-senal} verás, con los datos del curso, la primera gran
lección del machine learning: \textbf{sin señal no hay modelo}.

% ============================================================================
\section{¿Qué es machine learning para negocio?}
\label{sec:m7-que-es}

Piensa en la gerente de una sucursal con 15 años de experiencia. Con solo ver a
un cliente ``sabe'' si va a volver: se fija en cuánto tiempo lleva comprando, si
se ha quejado, si viene seguido. Nadie le dio una fórmula; la aprendió viendo
miles de casos. El machine learning hace lo mismo, pero con datos: revisa miles
de casos históricos, encuentra el patrón y lo aplica a casos nuevos, de forma
consistente y a escala (a 100{,}000 clientes en segundos).

\begin{concepto}{Machine learning}
El \termino{machine learning} es un conjunto de técnicas con las que la
computadora \textbf{aprende patrones a partir de datos históricos} para hacer
predicciones o encontrar grupos, sin que alguien programe las reglas a mano.
\begin{itemize}
  \item Un \termino{modelo} es una ``receta'' que convierte datos de entrada
        (antigüedad, quejas, uso\ldots) en una salida (probabilidad de
        abandono).
  \item \termino{Entrenar} un modelo es ajustar esa receta usando datos
        históricos donde ya conocemos la respuesta.
  \item \termino{Predecir} (o \emph{puntuar}) es aplicar la receta a datos
        nuevos, donde todavía no conocemos la respuesta.
\end{itemize}
\end{concepto}

\subsection{Supervisado y no supervisado}

Hay dos grandes familias de problemas, y distinguirlas es el primer paso para
elegir la herramienta correcta:

\begin{itemize}
  \item \termino{Aprendizaje supervisado}: tienes ejemplos históricos
        \emph{con la respuesta} (este cliente se fue, este no; esta tienda
        vendió \$850{,}000). El modelo aprende a predecir esa respuesta.
        \begin{itemize}
          \item \termino{Regresión}: la respuesta es un \textbf{número}
                (ventas, demanda, precio, días de retraso).
          \item \termino{Clasificación}: la respuesta es una \textbf{categoría}
                (se va / se queda; fraude / no fraude; compra / no compra).
        \end{itemize}
  \item \termino{Aprendizaje no supervisado}: \emph{no} hay una respuesta que
        predecir; buscas estructura oculta en los datos.
        \begin{itemize}
          \item \termino{Clustering} (agrupamiento): formar grupos de clientes,
                tiendas o productos parecidos entre sí.
          \item \termino{Reglas de asociación}: descubrir qué cosas ocurren
                juntas (quien compra laptop también compra mouse).
        \end{itemize}
\end{itemize}

Los \termino{pronósticos de series de tiempo} son un caso especial de
regresión: predicen un número futuro usando la historia de la propia variable
(tendencia y estacionalidad). La tabla~\ref{tab:m7-problemas} traduce preguntas
de negocio típicas a tipo de modelo y algoritmo.

\begin{table}[H]
\centering
\small
\caption{De la pregunta de negocio al algoritmo.}
\label{tab:m7-problemas}
\begin{tabularx}{\textwidth}{>{\raggedright\arraybackslash}X
                              >{\raggedright\arraybackslash}p{3.1cm}
                              >{\raggedright\arraybackslash}p{4.4cm}}
\toprule
\textbf{Pregunta de negocio} & \textbf{Tipo de problema} & \textbf{Algoritmos típicos} \\
\midrule
¿Cuánto venderá cada tienda el próximo mes? & Regresión & Regresión lineal, random forest \\
¿Qué clientes cancelarán su suscripción? & Clasificación & Regresión logística, árbol, random forest \\
¿Esta transacción es fraudulenta? & Clasificación (desbalanceada) & Regresión logística, random forest \\
¿Cuántas unidades pediremos al proveedor en diciembre? & Serie de tiempo & Ingenuo estacional, ETS, ARIMA \\
¿Qué grupos de clientes tenemos? & Clustering & RFM + k-means, clustering jerárquico \\
¿Qué productos se compran juntos? & Reglas de asociación & Soporte--confianza--lift, Apriori \\
¿Qué sucursales se parecen entre sí? & Clustering & k-means con variables escaladas \\
¿Aprobamos este crédito a un cliente? & Clasificación & Regresión logística (fácil de explicar) \\
¿Cuánto gastará un cliente en su vida con nosotros? & Regresión & Regresión lineal, random forest \\
\bottomrule
\end{tabularx}
\end{table}

\subsection{El flujo de trabajo de un proyecto de machine learning}
\label{sec:m7-flujo}

Un proyecto de machine learning no es ``correr un algoritmo''. Es un ciclo que
empieza y termina en el negocio (figura~\ref{fig:m7-flujo}).

\begin{figure}[H]
\centering
\begin{tikzpicture}[
  paso/.style={rectangle, rounded corners=3pt, draw=primario, thick,
    fill=primario!6, text width=2.75cm, minimum height=1.3cm, align=center,
    font=\sffamily\small},
  pasomodelo/.style={paso, draw=secundario, fill=secundario!8},
  flecha/.style={-{Stealth[length=2.5mm]}, thick, color=gristexto},
  node distance=0.65cm
]
\node[paso] (p1) {\textbf{1. Pregunta}\\de negocio};
\node[paso, right=of p1] (p2) {\textbf{2. Datos}\\históricos};
\node[paso, right=of p2] (p3) {\textbf{3. Variables}\\(features)};
\node[paso, right=of p3] (p4) {\textbf{4. Separar}\\entrenamiento\\y prueba};
\node[pasomodelo, below=1.3cm of p4] (p5) {\textbf{5. Entrenar}\\el modelo};
\node[pasomodelo, left=of p5] (p6) {\textbf{6. Evaluar}\\con datos de prueba};
\node[paso, left=of p6] (p7) {\textbf{7. Desplegar}\\(usar en el día a día)};
\node[paso, left=of p7] (p8) {\textbf{8. Monitorear}\\y reentrenar};
\draw[flecha] (p1) -- (p2);
\draw[flecha] (p2) -- (p3);
\draw[flecha] (p3) -- (p4);
\draw[flecha] (p4) -- (p5);
\draw[flecha] (p5) -- (p6);
\draw[flecha] (p6) -- (p7);
\draw[flecha] (p7) -- (p8);
\draw[flecha, dashed, color=acento] (p8.west) -- ++(-0.45,0) |- (p1.west)
  node[pos=0.25, left, font=\sffamily\scriptsize, color=acento,
       align=center] {el negocio\\cambia};
\draw[flecha, dashed, color=acento] (p6.north) -- (p3.south)
  node[midway, left, font=\sffamily\scriptsize, color=acento,
       align=right] {¿no alcanza?\\mejora variables};
\end{tikzpicture}
\caption{Flujo de trabajo de machine learning. Las flechas punteadas indican
que el proceso es un ciclo, no una línea recta.}
\label{fig:m7-flujo}
\end{figure}

\begin{enumerate}
  \item \textbf{Pregunta de negocio.} ``¿Qué clientes se van a ir en los
        próximos 3 meses?'' es una buena pregunta; ``haz algo con IA'' no lo
        es. Define también qué decisión tomarás con la respuesta.
  \item \textbf{Datos históricos.} Casos pasados donde ya conoces el resultado
        (clientes que se fueron y que se quedaron).
  \item \textbf{Variables (\termino{features}).} Las características que el
        modelo usará: antigüedad, número de quejas, uso del servicio\ldots{}
        Aquí se gana o se pierde la mayor parte del valor.
  \item \textbf{Separar entrenamiento y prueba.} Guardas una parte de los datos
        que el modelo \emph{nunca} ve al aprender, para evaluarlo honestamente.
  \item \textbf{Entrenar} uno o varios algoritmos.
  \item \textbf{Evaluar} con métricas que tengan sentido para el negocio.
  \item \textbf{Desplegar}: poner el modelo a trabajar (un CSV semanal para el
        call center, una columna en el tablero, una API).
  \item \textbf{Monitorear y reentrenar}: los clientes, los precios y la
        competencia cambian; un modelo que no se vigila se degrada.
\end{enumerate}

\begin{negocio}{El 80\% del trabajo no es el algoritmo}
En proyectos reales, la mayor parte del tiempo se va en los pasos 1 a 4:
entender la pregunta, conseguir y limpiar datos, y construir buenas variables
(todo lo que aprendiste en los Módulos~\ref{cap:manipulacion} a~\ref{cap:eda}).
Entrenar el modelo son unas cuantas líneas de código. Un modelo sencillo con
buenas variables casi siempre le gana a un modelo sofisticado con variables
pobres.
\end{negocio}

\subsection{Paquetes que usaremos}

Este módulo usa, además del \paquete{tidyverse}, estos paquetes: \paquete{cluster}
y \paquete{factoextra} (clustering), \paquete{rpart} (árboles de decisión),
\paquete{randomForest} (bosques aleatorios), \paquete{caret} (marco común para
entrenar y comparar modelos) y \paquete{forecast} (series de tiempo). Si no los
tienes, instálalos una sola vez:

\begin{rnoejecutar}
install.packages(c("cluster", "factoextra", "rpart", "randomForest",
                   "caret", "forecast"))
\end{rnoejecutar}

Y cárgalos al inicio de tu script:

\begin{rcode}
library(tidyverse)     # dplyr, ggplot2, readr, tidyr, purrr...
library(cluster)       # silueta y otras medidas de clustering
library(factoextra)    # gráficas de clustering (devuelven ggplot)
library(rpart)         # árboles de decisión
library(randomForest)  # bosques aleatorios
library(caret)         # entrenar y comparar modelos con un mismo formato
library(forecast)      # series de tiempo: ETS, ARIMA, pronósticos

# Comillas simples normales en las salidas de R (en lugar de tipográficas)
options(useFancyQuotes = FALSE)

# Colores corporativos del libro (Módulo 3)
color_principal <- "#2a78d6"
color_resalte   <- "#eb6834"
color_gris      <- "#b5b3ad"
paleta_libro <- c("#2a78d6", "#eb6834", "#1baf7a", "#eda100",
                  "#e87ba4", "#008300", "#4a3aa7", "#e34948")
\end{rcode}

Al cargarlos verás algunos mensajes normales, no errores: \paquete{caret}
avisa que carga \paquete{lattice}, \paquete{forecast} menciona un método de
\paquete{quantmod}, y \paquete{randomForest} avisa que ``enmascara'' las
funciones \codigo{combine()} de \paquete{dplyr} y \codigo{margin()} de
\paquete{ggplot2}. Eso significa que, si necesitas esas dos funciones, debes
escribirlas con su paquete: \codigo{dplyr::combine()} o
\codigo{ggplot2::margin()}.

\subsection{Sin señal no hay modelo}
\label{sec:m7-sin-senal}

Antes de entrenar cualquier cosa, veamos la lección más importante del módulo.
Un modelo solo puede aprender un patrón \emph{si el patrón existe en los
datos}. Intentemos explicar el monto de cada compra de
\archivo{transacciones.csv} con la tienda, el método de pago y el día de la
semana, usando una regresión lineal (la viste en el Módulo~\ref{cap:eda}).

\begin{rcode}
# Transacciones del curso (1,000 compras de 2023)
transacciones <- read_csv("datasets/transacciones.csv",
                          show_col_types = FALSE)

# ¿El monto depende de la tienda, el método de pago o el día?
modelo_sin_senal <- lm(
  total_transaccion ~ factor(tienda_id) + metodo_pago + dia_semana,
  data = transacciones
)

# R²: proporción de la variación del monto que explica el modelo (0 a 1)
summary(modelo_sin_senal)$r.squared
summary(modelo_sin_senal)$adj.r.squared   # R² ajustado
\end{rcode}

\begin{rsalida}
[1] 0.01406052
[1] -0.00403011
\end{rsalida}

El modelo explica apenas 1.4\% de la variación del monto, y el
\termino{R² ajustado} (que penaliza agregar variables inútiles) es incluso
negativo: el modelo es \emph{peor} que simplemente predecir el promedio. No es
un error de código: el generador de datos eligió tienda, método de pago y monto
de forma independiente, al azar. No hay nada que aprender.

\begin{advertencia}
Cuidado con la ``señal'' falsa. Si agregas \codigo{cantidad} al modelo anterior,
el R² sube a 0.45. ¿Descubrimos algo? No: en este archivo
\codigo{total\_transaccion} se \emph{calcula} como \codigo{cantidad} $\times$
\codigo{precio\_venta}, así que es aritmética, no un patrón de comportamiento.
Usar como predictora una variable que es parte del cálculo de la respuesta es
una forma de \textbf{fuga de datos} (sección~\ref{sec:m7-fuga}).
\end{advertencia}

Por eso, en las secciones de predicción simularemos datos donde \emph{sí}
existe una relación (por ejemplo, las quejas aumentan la probabilidad de
abandono). En tu empresa, los datos reales tendrán señal (y ruido); tu trabajo
será encontrarla y medir cuánta hay.

\begin{hazlotu}
Ajusta \codigo{lm(total\_transaccion \textasciitilde{} cantidad, data = transacciones)}
y revisa su R². Luego piensa: si mañana quisieras predecir el monto de una
compra \emph{antes} de que ocurra, ¿conocerías la cantidad? ¿Te sirve ese
modelo?
\end{hazlotu}

% ============================================================================
\section{Conceptos clave antes de entrenar}
\label{sec:m7-conceptos}

\subsection{Variables predictoras y variable objetivo}

En aprendizaje supervisado cada fila es una \termino{observación} (un cliente,
una tienda, un mes) y las columnas se dividen en dos tipos:

\begin{itemize}
  \item La \termino{variable objetivo} (\emph{target}, respuesta, $y$): lo que
        quieres predecir. Ejemplo: \codigo{churn} (si el cliente se fue).
  \item Las \termino{variables predictoras} (\emph{features}, $X$): la
        información que usarás para predecir. Ejemplo: antigüedad, quejas,
        uso, plan contratado.
\end{itemize}

\begin{table}[H]
\centering
\small
\caption{Ejemplo de datos para predecir abandono: cuatro predictoras y una
variable objetivo.}
\label{tab:m7-xy}
\begin{tabular}{lrrrlc}
\toprule
\textbf{Cliente} & \textbf{Antigüedad} & \textbf{Quejas} & \textbf{Uso mensual}
  & \textbf{Plan} & \textbf{¿Se fue? ($y$)} \\
\midrule
C0001 & 36 & 0 & 11 & Premium  & No \\
C0002 & 4  & 3 & 5  & Básico   & Sí \\
C0003 & 18 & 1 & 8  & Estándar & No \\
\multicolumn{5}{c}{\emph{variables predictoras ($X$)}} & \emph{objetivo} \\
\bottomrule
\end{tabular}
\end{table}

La regla de oro: \textbf{una predictora solo es válida si la conocerás en el
momento de predecir}. Si quieres saber hoy quién se irá el próximo mes, no
puedes usar ``motivo de cancelación'', porque ese dato solo existe cuando el
cliente ya se fue.

\subsection{Entrenamiento y prueba}
\label{sec:m7-split}

Imagina que preparas a un vendedor con 100 casos resueltos y luego lo evalúas
con \emph{esos mismos} 100 casos: sacará 10, pero no sabrás si aprendió o solo
memorizó. Con los modelos pasa igual. Por eso separamos los datos:

\begin{itemize}
  \item \termino{Conjunto de entrenamiento} (\emph{train}, 70--80\%): con él
        aprende el modelo.
  \item \termino{Conjunto de prueba} (\emph{test}, 20--30\%): se guarda bajo
        llave y solo se usa al final para medir qué tan bien predice casos que
        nunca vio.
\end{itemize}

Veámoslo con un ejemplo que usaremos en las siguientes secciones: 60
sucursales, cada una con su inversión mensual en publicidad y sus ventas del
mes (ambas en miles de pesos). Lo simulamos con una relación conocida:
\emph{ventas = 200 + 6 $\times$ publicidad + ruido}.

\begin{rcode}
set.seed(123)   # fija el azar: cualquiera que corra esto obtiene lo mismo
sucursales <- tibble(
  sucursal   = 1:60,
  publicidad = round(runif(60, min = 5, max = 50), 1),   # miles de $
  ventas     = round(200 + 6 * publicidad + rnorm(60, 0, 25), 1)
)

# Elegimos al azar 45 filas (75 %) para entrenar
set.seed(123)
filas_ent <- sample(nrow(sucursales), size = 45)

suc_entrenamiento <- sucursales[filas_ent, ]    # 45 sucursales
suc_prueba        <- sucursales[-filas_ent, ]   # las 15 restantes

nrow(suc_entrenamiento)
nrow(suc_prueba)
\end{rcode}

\begin{rsalida}
[1] 45
[1] 15
\end{rsalida}

\codigo{sample()} desordena al azar los números de fila y toma 45; el signo
menos en \codigo{sucursales[-filas\_ent, ]} significa ``todas las filas
\emph{excepto} esas''. Gracias a \codigo{set.seed(123)}, cada vez que ejecutes el
código obtendrás exactamente la misma división, algo indispensable para que tus
resultados sean reproducibles (y para que coincidan con los del libro).

\begin{consejo}
Para problemas de clasificación con una clase poco frecuente (por ejemplo, 20\%
de clientes que se van), usa \codigo{caret::createDataPartition()}: hace una
división \termino{estratificada}, es decir, mantiene la misma proporción de cada
clase en entrenamiento y prueba. Lo usaremos en la sección de churn.
\end{consejo}

\subsection{Sobreajuste: cuando el modelo memoriza}
\label{sec:m7-sobreajuste}

El \termino{sobreajuste} (\emph{overfitting}) ocurre cuando un modelo es tan
flexible que aprende de memoria los datos de entrenamiento, incluido su ruido,
y falla con datos nuevos. Es como el estudiante que memoriza las respuestas del
examen de práctica y reprueba el examen real.

Para verlo, supongamos que solo tuviéramos 14 sucursales para entrenar y
ajustemos polinomios cada vez más flexibles (grado 1 es una recta; grado 10 es
una curva con muchas ondas). Medimos el error con el \termino{RMSE} (raíz del
error cuadrático medio: el tamaño típico del error, en las mismas unidades que
las ventas; lo verás a detalle en la sección~\ref{sec:m7-metricas-reg}).

\begin{rcode}
# Función auxiliar: error típico de una predicción
rmse <- function(real, predicho) sqrt(mean((real - predicho)^2))

pocas_sucursales <- sucursales %>% slice(1:14)    # entrenamiento chico
resto_sucursales <- sucursales %>% slice(15:60)   # 46 casos "nuevos"

errores_por_grado <- map_dfr(1:10, function(g) {
  modelo <- lm(ventas ~ poly(publicidad, g), data = pocas_sucursales)
  tibble(
    grado = g,
    error_entrenamiento = rmse(pocas_sucursales$ventas, fitted(modelo)),
    error_prueba = rmse(resto_sucursales$ventas,
                        predict(modelo, resto_sucursales))
  )
})

errores_por_grado %>% mutate(across(-grado, ~ round(.x, 1)))
\end{rcode}

\begin{rsalida}
# A tibble: 10 x 3
   grado error_entrenamiento error_prueba
   <int>               <dbl>        <dbl>
 1     1                20.3         22  
 2     2                20.3         22.2
 3     3                20.2         22.1
 4     4                20.2         22  
 5     5                19.8         26.3
 6     6                17.6         56.6
 7     7                17.4         35.6
 8     8                15.9        669. 
 9     9                13.1       2945. 
10    10                10        14455. 
\end{rsalida}

Lee la tabla por columnas. El \textbf{error de entrenamiento} baja siempre que
el modelo se vuelve más flexible: la curva se acerca cada vez más a los 14
puntos. Pero el \textbf{error de prueba}, el que importa, empieza a subir en el
grado 5 y se dispara a partir del grado 8 (¡más de 14{,}000 en el grado 10!):
el modelo memorizó el ruido. La recta (grado 1), que coincide con
la forma real de los datos, es la que mejor generaliza. La
figura~\ref{fig:m07-sobreajuste} lo muestra visualmente.

\begin{rcode}
malla <- tibble(publicidad = seq(min(pocas_sucursales$publicidad),
                                 max(pocas_sucursales$publicidad),
                                 length.out = 300))
recta   <- lm(ventas ~ publicidad, data = pocas_sucursales)
memoria <- lm(ventas ~ poly(publicidad, 10), data = pocas_sucursales)

curvas <- bind_rows(
  malla %>% mutate(modelo = "Grado 1: generaliza",
                   ventas = predict(recta, malla)),
  malla %>% mutate(modelo = "Grado 10: memoriza",
                   ventas = predict(memoria, malla))
)

grafico_sobreajuste <- ggplot() +
  geom_point(data = resto_sucursales, aes(publicidad, ventas),
             color = color_gris, alpha = 0.8) +
  geom_point(data = pocas_sucursales, aes(publicidad, ventas),
             color = "black", size = 2.2) +
  geom_line(data = curvas, aes(publicidad, ventas, color = modelo),
            linewidth = 0.8) +
  scale_color_manual(values = c(color_principal, color_resalte)) +
  coord_cartesian(ylim = c(150, 600)) +
  labs(title = "Un modelo demasiado flexible memoriza el ruido",
       subtitle = paste("Negro: 14 sucursales de entrenamiento;",
                        "gris: sucursales nuevas"),
       x = "Publicidad (miles de $)", y = "Ventas (miles de $)",
       color = NULL) +
  theme_minimal() +
  theme(panel.grid.minor = element_blank(), legend.position = "top")
grafico_sobreajuste
\end{rcode}

\begin{figure}[H]
\centering
\includegraphics[width=0.9\textwidth]{m07-sobreajuste}
\caption{La recta (azul) captura la relación real; el polinomio de grado 10
(naranja) pasa cerca de cada punto de entrenamiento, pero hace predicciones
absurdas para sucursales nuevas.}
\label{fig:m07-sobreajuste}
\end{figure}

\begin{concepto}{Sesgo y varianza, en una frase}
Un modelo demasiado \textbf{simple} no captura el patrón (\emph{subajuste}); uno
demasiado \textbf{complejo} captura hasta el ruido (\emph{sobreajuste}). El
objetivo es el punto medio, y la única forma honesta de encontrarlo es medir el
error en datos que el modelo no usó para aprender.
\end{concepto}

\subsection{Validación cruzada}
\label{sec:m7-cv}

Con una sola división entrenamiento/prueba, el resultado depende un poco de la
suerte: quizá justo quedaron casos fáciles en la prueba. La
\termino{validación cruzada} de $k$ pliegues (\emph{k-fold cross-validation})
lo resuelve: parte los datos en $k$ grupos (típicamente 5 o 10), entrena $k$
veces dejando fuera un grupo distinto cada vez, y promedia los $k$ errores
(figura~\ref{fig:m7-cv}).

\begin{figure}[H]
\centering
\begin{tikzpicture}[
  bloque/.style={rectangle, draw=white, thick, minimum width=1.9cm,
    minimum height=0.55cm, font=\sffamily\scriptsize},
  ent/.style={bloque, fill=primario!25},
  val/.style={bloque, fill=acento!80, text=white},
]
\foreach \r in {1,...,5} {
  \node[font=\sffamily\small, anchor=east] at (-0.2, -\r*0.7) {Ronda \r};
  \foreach \c in {1,...,5} {
    \ifnum\c=\r
      \node[val] at (\c*1.95, -\r*0.7) {validación};
    \else
      \node[ent] at (\c*1.95, -\r*0.7) {entrena};
    \fi
  }
  \node[font=\sffamily\small, anchor=west] at (10.9, -\r*0.7) {error$_{\r}$};
}
\node[font=\sffamily\small, anchor=west, text=primario] at (10.9, -4.35)
  {\textbf{promedio}};
\end{tikzpicture}
\caption{Validación cruzada de 5 pliegues: cada sucursal se usa 4 veces para
entrenar y exactamente una vez para validar.}
\label{fig:m7-cv}
\end{figure}

Hagámoslo a mano una vez, para entender lo que después hará \paquete{caret}
automáticamente. Comparamos la recta con polinomios de grado 3, 6 y 10 usando
las 60 sucursales:

\begin{rcode}
set.seed(123)
# Asigna cada sucursal a uno de 5 pliegues, al azar y en partes iguales
pliegue <- sample(rep(1:5, length.out = nrow(sucursales)))
table(pliegue)

error_cv <- function(grado) {
  errores <- map_dbl(1:5, function(k) {
    entrena <- sucursales[pliegue != k, ]   # 4 pliegues para aprender
    valida  <- sucursales[pliegue == k, ]   # 1 pliegue para evaluar
    modelo  <- lm(ventas ~ poly(publicidad, grado), data = entrena)
    rmse(valida$ventas, predict(modelo, valida))
  })
  mean(errores)   # promedio de los 5 errores
}

tibble(grado = c(1, 3, 6, 10)) %>%
  mutate(rmse_cv = map_dbl(grado, error_cv))
\end{rcode}

\begin{rsalida}
pliegue
 1  2  3  4  5 
12 12 12 12 12 
# A tibble: 4 x 2
  grado rmse_cv
  <dbl>   <dbl>
1     1    21.5
2     3    22.2
3     6    22.6
4    10    29.0
\end{rsalida}

Primero, \codigo{table(pliegue)} confirma que cada pliegue tiene 12
sucursales. Luego, la recta tiene el menor error de validación cruzada (21.5
mil pesos), y el error crece con la complejidad hasta 29.0 con el grado 10. La
validación cruzada llega a la misma conclusión que la figura, pero usando las
60 sucursales y sin depender de una sola división afortunada. En la práctica
la usarás para \textbf{elegir entre modelos o ajustar sus parámetros}, y
reservarás el conjunto de prueba para la evaluación final.

\subsection{Fuga de datos: el error más caro}
\label{sec:m7-fuga}

La \termino{fuga de datos} (\emph{data leakage}) ocurre cuando, al entrenar, el
modelo tiene acceso a información que \textbf{no estará disponible en el
momento real de predecir}. El resultado es un modelo espectacular en el
laboratorio y inútil en la operación. Es el error más común y más caro del
machine learning, porque no produce ningún mensaje de error: produce
\emph{optimismo}.

\begin{negocio}{Un modelo de abandono ``perfecto''}
Un equipo entrena un modelo de churn y obtiene 98\% de exactitud. Entre las
predictoras incluyó \codigo{dias\_desde\_ultima\_llamada\_a\_cancelaciones}.
Pero los clientes llaman al área de cancelaciones \emph{cuando ya decidieron
irse}: esa variable es una consecuencia del abandono, no una causa observable
con anticipación. En producción, el modelo avisa cuando el cliente ya está
cancelando, demasiado tarde para retenerlo.
\end{negocio}

Formas típicas de fuga en BI:

\begin{itemize}
  \item \textbf{Variables del futuro}: motivo de cancelación, fecha de baja,
        monto de reembolso, ``estatus final'' del pedido.
  \item \textbf{Variables calculadas con la respuesta}: como el
        \codigo{total\_transaccion} que ya contiene a \codigo{cantidad}.
  \item \textbf{Preprocesar con todos los datos}: escalar, imputar o calcular
        promedios usando también los datos de prueba antes de separarlos.
  \item \textbf{Series de tiempo mezcladas}: separar al azar meses de una serie
        (entrenas con marzo de 2025 para ``predecir'' febrero de 2025). En series
        de tiempo, la prueba siempre son los periodos más recientes.
\end{itemize}

En la sección de errores comunes (sección~\ref{sec:m7-errores}) verás con
código cómo una sola variable con fuga convierte un modelo mediocre en uno
aparentemente perfecto.

\begin{hazlotu}
Piensa en tu empresa: escribe tres variables que \emph{parecerían} excelentes
para predecir si un pedido llegará tarde, pero que en realidad solo se conocen
después de la entrega.
\end{hazlotu}

% ============================================================================
\section{Caso 1: segmentación de clientes con RFM y k-means}
\label{sec:m7-segmentacion}

\begin{ejemplo}{¿A quién le mandamos la campaña VIP?}
El área de marketing tiene presupuesto para tres campañas distintas y quiere
dejar de mandar el mismo correo a los 200 clientes. La pregunta es: \emph{¿qué
grupos de clientes con comportamiento de compra parecido tenemos y qué
estrategia conviene para cada uno?} Es un problema \textbf{no supervisado}: no
hay una respuesta correcta que predecir; queremos descubrir grupos. Usaremos
los datos reales del curso: \archivo{transacciones.csv} y
\archivo{clientes.csv}.
\end{ejemplo}

\subsection{Paso 1: calcular Recencia, Frecuencia y Monto}

\begin{concepto}{RFM}
El análisis \termino{RFM} resume el comportamiento de cada cliente en tres
números, muy usados en retail y marketing directo:
\begin{itemize}
  \item \textbf{Recencia (R):} días desde su última compra. \emph{Menos} es
        mejor: quien compró hace poco tiene más probabilidad de volver.
  \item \textbf{Frecuencia (F):} número de compras en el periodo. Más es mejor.
  \item \textbf{Monto (M):} dinero total gastado en el periodo. Más es mejor.
\end{itemize}
La recencia se mide respecto a una \termino{fecha de corte} fija (por ejemplo,
el día en que haces el análisis). Fíjala en el código; si usas
\codigo{Sys.Date()}, tus resultados cambiarán cada día.
\end{concepto}

Las transacciones cubren todo 2023, así que usaremos como fecha de corte el
1 de enero de 2024:

\begin{rcode}
clientes <- read_csv("datasets/clientes.csv", show_col_types = FALSE)

fecha_corte <- as.Date("2024-01-01")   # fecha fija del análisis

rfm <- transacciones %>%
  group_by(cliente_id) %>%
  summarise(
    ultima_compra = max(fecha),
    recencia   = as.numeric(fecha_corte - ultima_compra),  # días
    frecuencia = n(),                                      # compras
    monto      = sum(total_transaccion),                   # pesos
    .groups = "drop"
  )

rfm
\end{rcode}

\begin{rsalida}
# A tibble: 198 x 5
   cliente_id ultima_compra recencia frecuencia  monto
        <dbl> <date>           <dbl>      <int>  <dbl>
 1          1 2023-12-01          31          7  4859.
 2          2 2023-12-02          30          5  7950.
 3          3 2023-11-04          58          8  9562 
 4          4 2023-12-24           8          5  7759.
 5          5 2023-12-31           1         13 14680.
 6          6 2023-12-31           1          6  6087.
 7          7 2023-12-13          19          5  1552.
 8          8 2023-12-05          27          5  3555.
 9          9 2023-12-23           9          6  5071.
10         10 2023-07-09         176          4  2487.
# i 188 more rows
\end{rsalida}

Cada fila es ahora un cliente. El cliente 1, por ejemplo, compró por última vez
el 1 de diciembre (31 días antes del corte), hizo 7 compras y gastó \$4{,}859.
Hay 198 filas, pero \archivo{clientes.csv} tiene 200 clientes. ¿Quiénes faltan?

\begin{rcode}
# Clientes registrados que no aparecen en las transacciones
clientes %>%
  anti_join(rfm, by = "cliente_id") %>%
  select(cliente_id, nombre, fecha_registro)
\end{rcode}

\begin{rsalida}
# A tibble: 2 x 3
  cliente_id nombre            fecha_registro
       <dbl> <chr>             <date>        
1        171 Rosa Jiménez      2022-08-27    
2        199 Ricardo Hernández 2023-09-16    
\end{rsalida}

Dos clientes no compraron nada en 2023. No entran al clustering (no tienen
recencia ni frecuencia), pero no los olvides: en la práctica formarían su propio
grupo, ``registrados sin compra'', con una estrategia de activación.

\begin{rcode}
rfm %>%
  select(recencia, frecuencia, monto) %>%
  summary()
\end{rcode}

\begin{rsalida}
    recencia        frecuencia         monto        
 Min.   :  1.00   Min.   : 1.000   Min.   :  152.6  
 1st Qu.: 21.00   1st Qu.: 3.000   1st Qu.: 3101.4  
 Median : 48.00   Median : 5.000   Median : 4521.8  
 Mean   : 69.19   Mean   : 5.051   Mean   : 4848.4  
 3rd Qu.: 98.75   3rd Qu.: 6.000   3rd Qu.: 6183.0  
 Max.   :314.00   Max.   :13.000   Max.   :14680.1  
\end{rsalida}

Observa las escalas: la recencia va de 1 a 314 días, la frecuencia de 1 a 13
compras y el monto de 153 a 14{,}680 pesos. Esto es crucial para el siguiente
paso.

\subsection{Paso 2: escalar las variables}

\begin{concepto}{Por qué escalar antes de k-means}
\termino{k-means} agrupa clientes según la \emph{distancia} entre ellos. Si una
variable se mide en miles (monto) y otra en unidades (frecuencia), la distancia
la domina casi por completo la de números grandes: una diferencia de 2 compras
``pesaría'' lo mismo que una diferencia de 2 pesos. \termino{Escalar}
(estandarizar) convierte cada variable a media 0 y desviación estándar 1, para
que las tres pesen lo mismo. En R se hace con \codigo{scale()}.
\end{concepto}

\begin{rcode}
rfm_escalado <- rfm %>%
  select(recencia, frecuencia, monto) %>%
  scale()                     # (x - media) / desviación estándar

round(colMeans(rfm_escalado), 2)          # medias: todas 0
round(apply(rfm_escalado, 2, sd), 2)      # desviaciones: todas 1
head(round(rfm_escalado, 2), 3)
\end{rcode}

\begin{rsalida}
  recencia frecuencia      monto 
         0          0          0 
  recencia frecuencia      monto 
         1          1          1 
     recencia frecuencia monto
[1,]    -0.59       0.85  0.00
[2,]    -0.60      -0.02  1.15
[3,]    -0.17       1.28  1.75
\end{rsalida}

Un valor escalado se lee como ``cuántas desviaciones estándar está este cliente
por arriba o por debajo del promedio''. El primer cliente tiene una recencia de
$-0.59$: compró más recientemente que el cliente promedio. \codigo{scale()}
guarda las medias y desviaciones usadas como atributos del resultado; los
necesitaremos después para asignar clientes nuevos.

\subsection{Paso 3: ¿cuántos grupos? Método del codo y silueta}

k-means necesita que le digas cuántos grupos ($k$) formar. Dos herramientas
ayudan a decidir:

\begin{itemize}
  \item \textbf{Método del codo:} para cada $k$ calcula la
        \termino{suma de cuadrados dentro de los grupos} (WSS, qué tan
        dispersos están los clientes respecto al centro de su grupo). Siempre
        baja al aumentar $k$; buscamos el ``codo'', el punto a partir del cual
        agregar grupos ya casi no mejora.
  \item \textbf{Silueta:} mide, para cada cliente, qué tan cerca está de su
        propio grupo comparado con el grupo vecino (de $-1$ a $1$; más alto es
        mejor). Buscamos el $k$ con la silueta promedio más alta.
\end{itemize}

\begin{rcode}
set.seed(123)
codo <- tibble(k = 1:10) %>%
  mutate(wss = map_dbl(k, function(k) {
    kmeans(rfm_escalado, centers = k, nstart = 25)$tot.withinss
  }))

codo %>%
  mutate(mejora_pct = round(100 * (lag(wss) - wss) / lag(wss), 1),
         wss = round(wss, 1))
\end{rcode}

\begin{rsalida}
# A tibble: 10 x 3
       k   wss mejora_pct
   <int> <dbl>      <dbl>
 1     1 591         NA  
 2     2 344.        41.9
 3     3 239         30.5
 4     4 175.        26.9
 5     5 146.        16.3
 6     6 127.        13  
 7     7 114         10.5
 8     8 102.        10.3
 9     9  92.3        9.7
10    10  85.1        7.8
\end{rsalida}

La columna \codigo{mejora\_pct} dice cuánto baja la dispersión al agregar un
grupo más. Pasar de 1 a 2 grupos reduce la WSS 41.9\%; de 2 a 3, 30.5\%; de 3 a
4, 26.9\%; y a partir de ahí las mejoras son cada vez menores. Veámoslo en una
gráfica, resaltando $k = 4$:

\begin{rcode}
grafico_codo <- ggplot(codo, aes(x = k, y = wss)) +
  geom_line(color = color_principal, linewidth = 0.8) +
  geom_point(color = color_principal, size = 2) +
  geom_point(data = filter(codo, k == 4), color = color_resalte,
             size = 4) +
  annotate("text", x = 4.3, y = codo$wss[4] + 60, hjust = 0,
           label = "k = 4: las mejoras\nempiezan a ser pequeñas",
           color = color_resalte, size = 3.5) +
  scale_x_continuous(breaks = 1:10) +
  labs(title = "Método del codo",
       subtitle = "Dispersión dentro de los grupos según k",
       x = "Número de grupos (k)",
       y = "Suma de cuadrados dentro (WSS)") +
  theme_minimal() +
  theme(panel.grid.minor = element_blank())
grafico_codo
\end{rcode}

\begin{figure}[H]
\centering
\includegraphics[width=0.9\textwidth]{m07-codo}
\caption{Método del codo para el RFM de los clientes del curso. La curva se
suaviza entre $k = 3$ y $k = 4$.}
\label{fig:m07-codo}
\end{figure}

Ahora la silueta. La función \codigo{fviz\_nbclust()} de \paquete{factoextra}
hace todo el cálculo y devuelve una gráfica de \paquete{ggplot2}, que puedes
personalizar con \codigo{labs()} y \codigo{theme()} como cualquier otra:

\begin{rcode}
set.seed(123)
grafico_silueta <- fviz_nbclust(rfm_escalado, kmeans,
                                method = "silhouette", k.max = 10,
                                nstart = 25, linecolor = color_principal) +
  labs(title = "Silueta promedio según el número de grupos",
       subtitle = "Más alto = grupos más separados",
       x = "Número de grupos (k)", y = "Silueta promedio") +
  theme_minimal() +
  theme(panel.grid.minor = element_blank())
grafico_silueta

# Los valores calculados vienen dentro del objeto ggplot
grafico_silueta$data %>%
  mutate(y = round(y, 3)) %>%
  filter(clusters %in% 2:6)
\end{rcode}

\begin{rsalida}
  clusters     y
1        2 0.345
2        3 0.337
3        4 0.328
4        5 0.325
5        6 0.292
\end{rsalida}

\begin{figure}[H]
\centering
\includegraphics[width=0.9\textwidth]{m07-silueta}
\caption{Silueta promedio para $k$ de 1 a 10. La línea punteada marca el
máximo ($k = 2$), aunque las diferencias entre 2 y 5 grupos son mínimas.}
\label{fig:m07-silueta}
\end{figure}

La silueta más alta es para $k = 2$ (0.345), pero $k = 3$, 4 y 5 están
prácticamente empatados (0.337, 0.328 y 0.325). Una silueta de alrededor de
0.33 indica una estructura \emph{débil}: los clientes no forman ``islas''
naturales, sino un continuo que k-means parte en regiones. Es lo esperable con
datos generados al azar, y muy común con clientes reales.

\begin{consejo}
Cuando las métricas no dan un ganador claro, decide con criterio de negocio.
Dos grupos (``buenos'' y ``malos'') no alcanzan para diseñar tres campañas;
diez grupos son imposibles de operar. Aquí elegimos $k = 4$: el codo lo
respalda, la silueta casi no empeora y el área de marketing puede manejar
cuatro estrategias. Las métricas orientan; la decisión es tuya.
\end{consejo}

\subsection{Paso 4: entrenar k-means}

\begin{concepto}{Cómo trabaja k-means}
\begin{enumerate}
  \item Coloca $k$ centros al azar.
  \item Asigna cada cliente al centro más cercano.
  \item Mueve cada centro al promedio de sus clientes.
  \item Repite 2 y 3 hasta que los grupos ya no cambian.
\end{enumerate}
Como el resultado depende de los centros iniciales, el argumento
\codigo{nstart = 25} prueba 25 arranques distintos y se queda con el mejor
(el de menor WSS). Usa siempre \codigo{set.seed()} y un \codigo{nstart} de 10
o más.
\end{concepto}

\begin{rcode}
set.seed(123)
modelo_kmeans <- kmeans(rfm_escalado, centers = 4, nstart = 25)

modelo_kmeans$size                     # clientes por grupo
round(modelo_kmeans$centers, 2)        # centros (en escala estandarizada)

# Proporción de la variación total explicada por los grupos
round(modelo_kmeans$betweenss / modelo_kmeans$totss, 3)
\end{rcode}

\begin{rsalida}
[1] 22 87 60 29
  recencia frecuencia monto
1    -0.53       1.81  1.97
2    -0.42       0.36  0.27
3    -0.13      -0.71 -0.73
4     1.95      -0.98 -0.80
[1] 0.704
\end{rsalida}

Los centros están en escala estandarizada, así que se leen como ``desviaciones
respecto al promedio''. El grupo 1 tiene frecuencia $+1.81$ y monto $+1.97$:
compra mucho más que el promedio. El grupo 4 tiene recencia $+1.95$: hace mucho
que no compra. Los cuatro grupos explican el 70.4\% de la variación total de los
datos.

\begin{advertencia}
Los números de grupo (1, 2, 3, 4) son etiquetas \textbf{arbitrarias}: si
cambias la semilla, el orden de los datos o la versión de R, el grupo de
``mejores clientes'' puede aparecer como 3 en lugar de 1. Nunca escribas en tu
código \codigo{cluster == 1 \textasciitilde{} "VIP"}. Nombra los grupos a partir
de su perfil calculado, como haremos enseguida.
\end{advertencia}

Revisemos también la calidad de cada grupo con la silueta individual:

\begin{rcode}
silueta <- silhouette(modelo_kmeans$cluster, dist(rfm_escalado))

# Silueta promedio por grupo y cuántos clientes quedaron "mal ubicados"
round(summary(silueta)$clus.avg.widths, 3)
sum(silueta[, "sil_width"] < 0)
\end{rcode}

\begin{rsalida}
    1     2     3     4 
0.375 0.321 0.298 0.379 
[1] 0
\end{rsalida}

Todos los grupos tienen silueta positiva; los más compactos son el 4 (0.379)
y el 1 (0.375), y el menos definido es el 3 (0.298). Ningún cliente tiene
silueta negativa (el \codigo{0} de la última línea): nadie quedó más cerca del
grupo vecino que del suyo. Si hubiera clientes con silueta negativa, serían
clientes ``frontera'' que conviene revisar a mano.

\subsection{Paso 5: perfilar y nombrar los segmentos}

Un segmento no sirve hasta que alguien de negocio puede describirlo en una
frase. Calculamos el perfil de cada grupo \emph{en las unidades originales}
(días, compras, pesos) y su peso en clientes e ingresos:

\begin{rcode}
rfm_seg <- rfm %>%
  mutate(cluster = modelo_kmeans$cluster)

perfil <- rfm_seg %>%
  group_by(cluster) %>%
  summarise(
    clientes        = n(),
    recencia_prom   = mean(recencia),
    frecuencia_prom = mean(frecuencia),
    monto_prom      = mean(monto),
    ingresos        = sum(monto),
    .groups = "drop"
  ) %>%
  mutate(pct_clientes = 100 * clientes / sum(clientes),
         pct_ingresos = 100 * ingresos / sum(ingresos))

perfil %>%
  select(cluster, clientes, recencia_prom, frecuencia_prom, monto_prom,
         pct_ingresos) %>%
  mutate(across(where(is.double), ~ round(.x, 1)))
\end{rcode}

\begin{rsalida}
# A tibble: 4 x 6
  cluster clientes recencia_prom frecuencia_prom monto_prom
    <int>    <int>         <dbl>           <dbl>      <dbl>
1       1       22          34.5             9.2     10159.
2       2       87          41.8             5.9      5577.
3       3       60          60.5             3.4      2879.
4       4       29         196.              2.8      2708.
# i 1 more variable: pct_ingresos <dbl>
\end{rsalida}

Ahora asignamos los nombres \textbf{con una regla basada en el perfil}, no en
el número del grupo. La regla de negocio es: el grupo con la recencia promedio
más alta son los ``Dormidos''; los demás se ordenan por monto promedio de mayor
a menor y se llaman ``Campeones'', ``Leales'' y ``Ocasionales''.

\begin{rcode}
nombres_segmentos <- perfil %>%
  mutate(es_dormido = recencia_prom == max(recencia_prom)) %>%
  arrange(es_dormido, desc(monto_prom)) %>%   # dormidos al final
  mutate(segmento = c("Campeones", "Leales", "Ocasionales",
                      "Dormidos")) %>%
  select(cluster, segmento)

orden_segmentos <- c("Campeones", "Leales", "Ocasionales", "Dormidos")

rfm_seg <- rfm_seg %>%
  left_join(nombres_segmentos, by = "cluster") %>%
  mutate(segmento = factor(segmento, levels = orden_segmentos))

perfil_segmentos <- perfil %>%
  left_join(nombres_segmentos, by = "cluster") %>%
  mutate(segmento = factor(segmento, levels = orden_segmentos)) %>%
  arrange(segmento) %>%
  transmute(segmento, clientes,
            recencia_dias = round(recencia_prom),
            compras       = round(frecuencia_prom, 1),
            monto_prom    = round(monto_prom),
            pct_ingresos  = round(pct_ingresos, 1))
perfil_segmentos
\end{rcode}

\begin{rsalida}
# A tibble: 4 x 6
  segmento    clientes recencia_dias compras monto_prom pct_ingresos
  <fct>          <int>         <dbl>   <dbl>      <dbl>        <dbl>
1 Campeones         22            35     9.2      10159         23.3
2 Leales            87            42     5.9       5577         50.5
3 Ocasionales       60            60     3.4       2879         18  
4 Dormidos          29           195     2.8       2708          8.2
\end{rsalida}

Ahora sí, cada grupo cuenta una historia:

\begin{itemize}
  \item \textbf{Campeones} (22 clientes, 11\% de la base): compraron hace unos
        35 días, 9.2 veces en el año y gastaron en promedio \$10{,}159. Son
        apenas uno de cada nueve clientes, pero generan el 23.3\% de los
        ingresos.
  \item \textbf{Leales} (87 clientes): el grueso de la base. Compran seguido
        (5.9 veces) y aportan la mitad de los ingresos (50.5\%).
  \item \textbf{Ocasionales} (60): activos (60 días de recencia), pero compran
        poco y gastan \$2{,}879 en promedio.
  \item \textbf{Dormidos} (29): su última compra fue hace 195 días en
        promedio, más de seis meses. Aportan solo 8.2\% de los ingresos.
\end{itemize}

Observa que la regla de nombres no menciona ningún número de grupo: si mañana
k-means numera distinto, cada cliente seguirá recibiendo el nombre correcto.

\begin{hazlotu}
Cambia la semilla a \codigo{set.seed(2024)} antes de \codigo{kmeans()} y vuelve
a correr los pasos 4 y 5. ¿Cambian los números de grupo? ¿Cambian los nombres
asignados a cada cliente? (Deberían mantenerse, porque los nombres dependen del
perfil y no del número.)
\end{hazlotu}

\subsection{Paso 6: visualizar los segmentos}

\codigo{fviz\_cluster()} dibuja los grupos en dos dimensiones. Como tenemos tres
variables, primero las resume en dos ``componentes principales'' (combinaciones
de R, F y M que conservan la mayor parte de la información); los porcentajes de
los ejes indican cuánta información conserva cada una.

\begin{rcode}
grafico_clusters <- fviz_cluster(modelo_kmeans, data = rfm_escalado,
                                 geom = "point", ellipse.type = "convex",
                                 palette = paleta_libro[1:4],
                                 ggtheme = theme_minimal()) +
  labs(title = "Grupos de k-means sobre el RFM escalado",
       subtitle = "Proyección en los dos componentes principales",
       color = "Grupo", fill = "Grupo", shape = "Grupo")
grafico_clusters
\end{rcode}

\begin{figure}[H]
\centering
\includegraphics[width=0.9\textwidth]{m07-clusters}
\caption{Los cuatro grupos de k-means proyectados en dos componentes
principales, que juntos conservan 93.2\% de la información de R, F y M.}
\label{fig:m07-clusters}
\end{figure}

La gráfica de componentes principales es útil para ver si los grupos se
separan, pero sus ejes son difíciles de explicar a la dirección. Para una
presentación, grafica las variables originales con los nombres de negocio:

\begin{rcode}
grafico_segmentos <- ggplot(rfm_seg,
                            aes(recencia, monto, color = segmento)) +
  geom_point(size = 2, alpha = 0.8) +
  scale_color_manual(values = paleta_libro[1:4]) +
  scale_y_continuous(labels = scales::dollar) +
  labs(title = "Segmentos de clientes según recencia y monto",
       subtitle = "Clientes del curso, compras de 2023",
       x = "Días desde la última compra (recencia)",
       y = "Monto total comprado", color = NULL) +
  theme_minimal() +
  theme(panel.grid.minor = element_blank(), legend.position = "top")
grafico_segmentos
\end{rcode}

\begin{figure}[H]
\centering
\includegraphics[width=0.9\textwidth]{m07-segmentos}
\caption{Segmentos RFM en unidades de negocio: los Dormidos se separan por
recencia; Campeones, Leales y Ocasionales, por monto.}
\label{fig:m07-segmentos}
\end{figure}

\subsection{Paso 7: de segmentos a estrategias}

La segmentación solo genera valor cuando cada segmento tiene una acción
asociada. La tabla~\ref{tab:m7-estrategias} es un ejemplo de cómo se
presentaría a marketing.

\begin{table}[H]
\centering
\small
\caption{Estrategia sugerida por segmento RFM.}
\label{tab:m7-estrategias}
\begin{tabularx}{\textwidth}{>{\raggedright\arraybackslash}p{2.2cm}
                              >{\raggedright\arraybackslash}p{4.2cm}
                              >{\raggedright\arraybackslash}X}
\toprule
\textbf{Segmento} & \textbf{Quiénes son} & \textbf{Estrategia} \\
\midrule
Campeones & Compran seguido, gastan el doble del promedio y compraron hace poco. &
  Programa VIP, acceso anticipado a lanzamientos, atención preferente. No
  regalar descuentos: ya compran sin ellos. \\
Leales & Compradores regulares con gasto por encima del promedio. &
  Venta cruzada y \emph{up-selling} (sección~\ref{sec:m7-canasta}); programa
  de puntos para convertirlos en Campeones. \\
Ocasionales & Compran poco y gastan poco, pero siguen activos. &
  Campañas de frecuencia: cupones con fecha de vencimiento, recordatorios,
  paquetes de entrada. \\
Dormidos & No compran desde hace más de seis meses en promedio. &
  Campaña de reactivación con oferta fuerte y encuesta de salida; si no
  responden, reducir la inversión en ellos. \\
\bottomrule
\end{tabularx}
\end{table}

Para enriquecer el perfil, puedes unir los segmentos con los datos
demográficos de \archivo{clientes.csv}. Por ejemplo, ¿los Campeones son más
``premium'' o más jóvenes?

\begin{rcode}
rfm_seg %>%
  # de clientes.csv solo tomamos lo que necesitamos (ya trae una
  # columna "segmento" por edad que chocaría con la nuestra)
  left_join(select(clientes, cliente_id, edad, es_premium),
            by = "cliente_id") %>%
  group_by(segmento) %>%
  summarise(clientes    = n(),
            edad_prom   = round(mean(edad), 1),
            pct_premium = round(100 * mean(es_premium), 1),
            .groups = "drop")
\end{rcode}

\begin{rsalida}
# A tibble: 4 x 4
  segmento    clientes edad_prom pct_premium
  <fct>          <int>     <dbl>       <dbl>
1 Campeones         22      49.2        31.8
2 Leales            87      47.4        34.5
3 Ocasionales       60      49.2        45  
4 Dormidos          29      51.3        27.6
\end{rsalida}

La edad promedio es casi igual en los cuatro segmentos (47 a 51 años) y el
porcentaje de clientes \emph{premium} no sigue ningún patrón lógico: los
Ocasionales tienen más premium (45\%) que los Campeones (31.8\%). En estos datos
simulados, la membresía premium y la edad no tienen relación con el
comportamiento de compra. En una empresa real este cruce sería un hallazgo
importante: si pagas beneficios premium a clientes que compran poco, el
programa necesita revisión.

\subsection{Paso 8: asignar un segmento a clientes nuevos}

El modelo ya está entrenado. Cuando un cliente acumula historial, no necesitas
volver a correr k-means: basta con escalar sus valores \textbf{con las mismas
medias y desviaciones del entrenamiento} y asignarlo al centro más cercano.

\begin{rcode}
asignar_segmento <- function(nuevos, modelo, datos_escalados, nombres) {
  medias   <- attr(datos_escalados, "scaled:center")
  desv_est <- attr(datos_escalados, "scaled:scale")
  # Escalamos con los parámetros del entrenamiento (no con los nuevos)
  nuevos_esc <- scale(nuevos[, names(medias)], center = medias,
                      scale = desv_est)
  # Para cada cliente: el centro con menor distancia al cuadrado
  grupo <- apply(nuevos_esc, 1, function(fila) {
    which.min(colSums((t(modelo$centers) - fila)^2))
  })
  nuevos %>%
    mutate(cluster = grupo) %>%
    left_join(nombres, by = "cluster")
}

nuevos_clientes <- tibble(
  cliente    = c("Ana", "Beto", "Carla"),
  recencia   = c(10, 250, 45),
  frecuencia = c(11, 2, 4),
  monto      = c(9500, 1800, 3500)
)

asignar_segmento(nuevos_clientes, modelo_kmeans, rfm_escalado,
                 nombres_segmentos)
\end{rcode}

\begin{rsalida}
# A tibble: 3 x 6
  cliente recencia frecuencia monto cluster segmento   
  <chr>      <dbl>      <dbl> <dbl>   <int> <chr>      
1 Ana           10         11  9500       1 Campeones  
2 Beto         250          2  1800       4 Dormidos   
3 Carla         45          4  3500       3 Ocasionales
\end{rsalida}

Ana (compró hace 10 días, 11 compras, \$9{,}500) es Campeona; Beto, que no
compra desde hace 250 días, cae en Dormidos; y Carla, con un comportamiento
modesto pero reciente, en Ocasionales. Este tipo de función es la que después
se programa para que corra cada semana (Módulo~\ref{cap:reportes}).

\subsection{Alternativa sencilla: RFM por reglas con quintiles}

No siempre necesitas k-means. Muchas empresas usan \termino{RFM por reglas}: a
cada cliente se le asigna una calificación de 1 a 5 en cada dimensión según su
quintil (el 20\% con mejor valor recibe 5) y luego se definen segmentos con
reglas fáciles de explicar. En R, \codigo{ntile()} de \paquete{dplyr} hace el
trabajo.

\begin{rcode}
rfm_reglas <- rfm %>%
  mutate(
    r_score = ntile(desc(recencia), 5),   # 5 = compró más recientemente
    f_score = ntile(frecuencia, 5),       # 5 = compra más seguido
    m_score = ntile(monto, 5),            # 5 = gasta más
    codigo_rfm = paste0(r_score, f_score, m_score),
    segmento_reglas = case_when(
      r_score >= 4 & f_score >= 4 & m_score >= 4 ~ "Campeones",
      r_score >= 3 & f_score >= 3                ~ "Leales",
      r_score <= 2 & f_score >= 3                ~ "En riesgo",
      r_score <= 2                               ~ "Dormidos",
      TRUE                                       ~ "Ocasionales"
    )
  )

rfm_reglas %>%
  select(cliente_id, recencia, frecuencia, monto, codigo_rfm,
         segmento_reglas) %>%
  head(5)

rfm_reglas %>% count(segmento_reglas, sort = TRUE)
\end{rcode}

\begin{rsalida}
# A tibble: 5 x 6
  cliente_id recencia frecuencia  monto codigo_rfm segmento_reglas
       <dbl>    <dbl>      <int>  <dbl> <chr>      <chr>          
1          1       31          7  4859. 443        Leales         
2          2       30          5  7950. 435        Leales         
3          3       58          8  9562  355        Leales         
4          4        8          5  7759. 535        Leales         
5          5        1         13 14680. 555        Campeones      
# A tibble: 5 x 2
  segmento_reglas     n
  <chr>           <int>
1 Leales             56
2 Dormidos           50
3 Campeones          32
4 En riesgo          30
5 Ocasionales        30
\end{rsalida}

El cliente 1 tiene código 443 (recencia y frecuencia altas, monto medio): no
alcanza el 4 en monto para ser Campeón y cae en ``Leales''. Ahora un segmento nuevo aparece: ``En riesgo'', clientes que
\emph{solían} comprar seguido pero llevan tiempo sin volver; k-means los había
repartido entre otros grupos. ¿Qué tanto coinciden ambos métodos?

\begin{rcode}
table(Reglas = rfm_reglas$segmento_reglas, kmeans = rfm_seg$segmento)
\end{rcode}

\begin{rsalida}
             kmeans
Reglas        Campeones Leales Ocasionales Dormidos
  Campeones          12     20           0        0
  Dormidos            0      1          23       26
  En riesgo           3     19           5        3
  Leales              7     44           5        0
  Ocasionales         0      3          27        0
\end{rsalida}

Las filas son los segmentos por reglas y las columnas los de k-means. La
coincidencia es buena en los extremos: 27 de los 30 Ocasionales por reglas son
Ocasionales en k-means y 26 de los 29 Dormidos de k-means son Dormidos por
reglas. Los 30 clientes ``En riesgo'' se reparten entre Leales (19) y otros
grupos: k-means los mezcló con clientes activos porque su monto y frecuencia
pesan más que su recencia. Ningún método es ``el correcto''; son dos formas de
mirar a los mismos clientes.

\begin{negocio}{¿Reglas o k-means?}
Las reglas RFM son transparentes (``Campeón = calificación 4 o 5 en todo''),
estables en el tiempo y fáciles de programar incluso en SQL o Excel. k-means
descubre grupos sin que tú fijes los cortes y admite más variables (margen,
categorías compradas, devoluciones), pero sus grupos pueden moverse al
reentrenar. Muchas empresas empiezan con reglas y pasan a clustering cuando
necesitan más variables.
\end{negocio}

\begin{hazlotu}
Agrega al RFM una cuarta variable: el número de tiendas distintas en las que
compró cada cliente (\codigo{n\_distinct(tienda\_id)}). Escala, vuelve a correr
k-means con $k = 4$ y compara los perfiles. ¿Aparece algún grupo nuevo con
sentido de negocio?
\end{hazlotu}

% ============================================================================
\section{Caso 2: regresión para predecir ventas}
\label{sec:m7-regresion}

\begin{ejemplo}{¿Cuánto venderá una tienda nueva?}
Una cadena con 400 tiendas planea abrir sucursales y quiere estimar sus ventas
mensuales antes de firmar el contrato de renta. Para cada tienda existente
conoce su zona, tamaño en m², tráfico peatonal mensual, inversión en
publicidad local y si participa en el programa de promociones. La pregunta:
\emph{¿cuánto venderá una tienda con ciertas características y qué palancas
mueven más las ventas?} Es un problema de \textbf{regresión}: la respuesta es
un número.
\end{ejemplo}

\subsection{Datos simulados con una relación conocida}

Como explicamos al inicio, los datos del curso no tienen esta información
(\archivo{tiendas.csv} solo tiene 10 tiendas y sin publicidad), así que
simulamos 400 tiendas. Las ventas mensuales (en miles de pesos) se generan con
esta ``verdad'' que el modelo tendrá que descubrir:

\[
\text{ventas} = 120 + 3.2\,\text{publicidad} + 0.25\,\text{m}^2 +
0.04\,\text{tráfico} + 60\,\text{promoción} + \text{efecto zona} + \text{ruido}
\]

con efecto zona de $+40$ en el Norte, $-30$ en el Sur y 0 en el Centro. El
tráfico depende en parte del tamaño (las tiendas grandes atraen más gente),
como en la vida real.

\begin{rcode}
set.seed(123)
n_tiendas <- 400
tiendas_sim <- tibble(
  tienda     = sprintf("T%03d", 1:n_tiendas),
  zona       = sample(c("Centro", "Norte", "Sur"), n_tiendas,
                      replace = TRUE),
  tamano_m2  = round(runif(n_tiendas, 150, 1200)),
  publicidad = round(runif(n_tiendas, 10, 120), 1),  # miles de $/mes
  promocion  = rbinom(n_tiendas, 1, 0.35),           # 1 = participa
  trafico    = round(3000 + 5 * tamano_m2 +          # personas/mes
                       rnorm(n_tiendas, 0, 1500))
) %>%
  mutate(
    efecto_zona = case_when(zona == "Norte" ~ 40,
                            zona == "Sur"   ~ -30,
                            TRUE            ~ 0),
    ventas = round(120 + 3.2 * publicidad + 0.25 * tamano_m2 +
                     0.04 * trafico + 60 * promocion + efecto_zona +
                     rnorm(n_tiendas, 0, 50), 1)       # miles de $/mes
  ) %>%
  select(-efecto_zona)   # el modelo no debe ver la "respuesta secreta"

head(tiendas_sim, 4)
\end{rcode}

\begin{rsalida}
# A tibble: 4 x 7
  tienda zona  tamano_m2 publicidad promocion trafico ventas
  <chr>  <chr>     <dbl>      <dbl>     <int>   <dbl>  <dbl>
1 T001   Sur         508       17.9         0    5980   517.
2 T002   Sur         638       10.4         0    3903   408.
3 T003   Sur         237       15.7         0    5900   400.
4 T004   Norte      1053      105.          1   10671  1263.
\end{rsalida}

Separamos 80\% para entrenar y 20\% para probar:

\begin{rcode}
set.seed(123)
filas_ent <- sample(nrow(tiendas_sim), size = 0.8 * nrow(tiendas_sim))
tiendas_ent    <- tiendas_sim[filas_ent, ]    # 320 tiendas
tiendas_prueba <- tiendas_sim[-filas_ent, ]   # 80 tiendas
c(entrenamiento = nrow(tiendas_ent), prueba = nrow(tiendas_prueba))
\end{rcode}

\begin{rsalida}
entrenamiento        prueba 
          320            80 
\end{rsalida}

\subsection{Regresión lineal múltiple}

La \termino{regresión lineal múltiple} (función \codigo{lm()}, que ya usaste en
el Módulo~\ref{cap:eda} con una sola variable) estima cuánto cambia la
respuesta por cada unidad de cada predictora, \emph{manteniendo las demás
constantes}. La fórmula \codigo{ventas \textasciitilde{} publicidad + tamano\_m2 + \ldots}
se lee ``ventas en función de publicidad, tamaño, \ldots''.

\begin{rcode}
modelo_lm <- lm(ventas ~ publicidad + tamano_m2 + trafico + promocion +
                  zona, data = tiendas_ent)
summary(modelo_lm)
\end{rcode}

\begin{rsalida}

Call:
lm(formula = ventas ~ publicidad + tamano_m2 + trafico + promocion + 
    zona, data = tiendas_ent)

Residuals:
     Min       1Q   Median       3Q      Max 
-120.181  -38.036   -1.857   35.175  141.550 

Coefficients:
              Estimate Std. Error t value Pr(>|t|)    
(Intercept) 107.296043  11.071131   9.692  < 2e-16 ***
publicidad    3.259854   0.091005  35.821  < 2e-16 ***
tamano_m2     0.270330   0.013935  19.399  < 2e-16 ***
trafico       0.038681   0.001903  20.330  < 2e-16 ***
promocion    57.226212   6.042430   9.471  < 2e-16 ***
zonaNorte    41.704523   6.880492   6.061 3.88e-09 ***
zonaSur     -29.067746   6.992714  -4.157 4.17e-05 ***
---
Signif. codes:  0 '***' 0.001 '**' 0.01 '*' 0.05 '.' 0.1 ' ' 1

Residual standard error: 50.65 on 313 degrees of freedom
Multiple R-squared:  0.937,	Adjusted R-squared:  0.9358 
F-statistic:   776 on 6 and 313 DF,  p-value: < 2.2e-16
\end{rsalida}

\begin{concepto}{Cómo leer la tabla de coeficientes}
\begin{itemize}
  \item \textbf{Estimate}: el efecto estimado. Para una variable numérica, es
        el cambio en ventas por cada unidad adicional; para una categoría
        (\codigo{zonaNorte}), es la diferencia respecto a la categoría base
        (Centro, la primera en orden alfabético).
  \item \textbf{Pr(>|t|)} (valor p): si es muy pequeño (menor a 0.05), el
        efecto difícilmente se debe al azar. Los asteriscos lo resumen.
  \item \textbf{Residual standard error}: el tamaño típico del error del
        modelo, en unidades de la respuesta.
  \item \textbf{R-squared}: proporción de la variación de las ventas que el
        modelo explica.
\end{itemize}
\end{concepto}

Traducido al lenguaje de negocio (recuerda que ventas y publicidad están en
miles de pesos):

\begin{itemize}
  \item \textbf{Publicidad:} cada \$1{,}000 adicionales de publicidad se asocian
        con \$3{,}260 más de ventas mensuales. El valor real era 3.2: el modelo
        lo descubrió.
  \item \textbf{Tamaño:} cada m² adicional aporta unos \$270 al mes (real: 250).
  \item \textbf{Tráfico:} cada 100 visitantes más, \$3{,}900 más de ventas
        (0.039 $\times$ 100, real: 0.04).
  \item \textbf{Promoción:} participar suma unos \$57{,}000 al mes (real: 60).
  \item \textbf{Zona:} una tienda del Norte vende \$41{,}700 más que una
        equivalente del Centro; una del Sur, \$29{,}100 menos.
  \item El modelo explica el 93.7\% de la variación de las ventas
        (R² = 0.937), y el error típico es de unos \$50{,}000, justo el tamaño
        del ruido que simulamos.
\end{itemize}

\begin{advertencia}
Un coeficiente de regresión mide \textbf{asociación, no causalidad}. En datos
simulados sabemos que la publicidad \emph{causa} ventas, porque así lo
programamos. En datos reales, quizá las tiendas que ya venden mucho reciben más
presupuesto de publicidad (la causalidad va al revés). Para afirmar ``si
invierto \$1{,}000 más en publicidad venderé \$3{,}260 más'' necesitas un
experimento (por ejemplo, aumentar la publicidad en tiendas elegidas al azar).
\end{advertencia}

\begin{consejo}
Para saber qué variable ``pesa'' más, no compares los coeficientes
directamente: están en unidades distintas (pesos, m², personas). Compara el
efecto de un cambio \emph{realista}: pasar de 40 a 80 mil pesos de publicidad
suma $40 \times 3.26 \approx 130$ mil pesos; pasar de 400 a 800 m² suma
$400 \times 0.27 \approx 108$ mil pesos directamente, más el tráfico extra que
atrae.
\end{consejo}

\subsection{Evaluar el modelo: MAE, RMSE, MAPE y R²}
\label{sec:m7-metricas-reg}

El resumen del modelo describe el ajuste en los datos de entrenamiento. La
evaluación honesta se hace con las 80 tiendas de prueba. Las cuatro métricas
de regresión más usadas son:

\begin{table}[H]
\centering
\small
\caption{Métricas de error para regresión.}
\label{tab:m7-metricas-reg}
\begin{tabularx}{\textwidth}{>{\raggedright\arraybackslash}p{2.2cm}
                              >{\raggedright\arraybackslash}p{5.4cm}
                              >{\raggedright\arraybackslash}X}
\toprule
\textbf{Métrica} & \textbf{Cálculo} & \textbf{Cómo se lee} \\
\midrule
MAE (error absoluto medio) & promedio de $|\text{real} - \text{predicho}|$ &
  ``En promedio nos equivocamos por \$X''. La más fácil de explicar. \\
RMSE (raíz del error cuadrático medio) &
  $\sqrt{\text{promedio de } (\text{real} - \text{predicho})^2}$ &
  Como el MAE, pero castiga más los errores grandes. Útil si un error grande
  es muy costoso. \\
MAPE (error porcentual absoluto medio) &
  promedio de $|\text{real} - \text{predicho}| / \text{real} \times 100$ &
  ``Nos equivocamos X\% en promedio''. Comparable entre negocios de distinto
  tamaño; falla si hay valores reales cercanos a cero. \\
R² (coeficiente de determinación) &
  $1 - \dfrac{\sum (\text{real} - \text{pred.})^2}{\sum (\text{real} - \text{prom.})^2}$ &
  Proporción de la variación explicada: 1 es perfecto, 0 es igual a predecir
  siempre el promedio, negativo es peor que el promedio. \\
\bottomrule
\end{tabularx}
\end{table}

Calculémoslas a mano con una función reutilizable:

\begin{rcode}
metricas_regresion <- function(real, predicho) {
  error <- real - predicho
  tibble(
    MAE  = mean(abs(error)),
    RMSE = sqrt(mean(error^2)),
    MAPE = mean(abs(error) / real) * 100,
    R2   = 1 - sum(error^2) / sum((real - mean(real))^2)
  )
}

tiendas_prueba <- tiendas_prueba %>%
  mutate(pred_lm = predict(modelo_lm, newdata = tiendas_prueba))

metricas_regresion(tiendas_prueba$ventas, tiendas_prueba$pred_lm) %>%
  mutate(across(everything(), ~ round(.x, 3)))
\end{rcode}

\begin{rsalida}
# A tibble: 1 x 4
    MAE  RMSE  MAPE    R2
  <dbl> <dbl> <dbl> <dbl>
1  37.2  45.1  5.04 0.947
\end{rsalida}

En tiendas que el modelo nunca vio, nos equivocamos en promedio por \$37{,}200
al mes (MAE = 37.2 miles), lo que equivale a un error de 5\% (MAPE). El RMSE
(45.1) es algo mayor que el MAE porque penaliza más los errores grandes; si
fuera mucho mayor, indicaría que hay algunas tiendas con errores enormes. El R²
en prueba (0.947) es tan bueno como en entrenamiento (0.937): no hay señales de
sobreajuste.

\begin{consejo}
Compara siempre contra una \termino{línea base} (\emph{baseline}): el modelo
más tonto posible. Aquí, predecir para toda tienda nueva el promedio de ventas
del entrenamiento. Si tu modelo no le gana claramente a la línea base, no
aporta valor.
\end{consejo}

\begin{rcode}
prediccion_ingenua <- mean(tiendas_ent$ventas)   # siempre el promedio
metricas_regresion(tiendas_prueba$ventas,
                   rep(prediccion_ingenua, nrow(tiendas_prueba))) %>%
  mutate(across(everything(), ~ round(.x, 3)))
\end{rcode}

\begin{rsalida}
# A tibble: 1 x 4
    MAE  RMSE  MAPE     R2
  <dbl> <dbl> <dbl>  <dbl>
1  154.  196.  24.7 -0.003
\end{rsalida}

Predecir siempre el promedio tiene un MAE de 154 mil pesos y un MAPE de
24.7\%; su R² es prácticamente cero (por definición, el promedio no explica
nada). Nuestro modelo reduce el error promedio en más de tres cuartas partes:
de 154 a 37 mil pesos.

\subsection{Gráfico real contra predicho}

La mejor forma de mostrar la calidad de un modelo de regresión a la dirección
es un diagrama de dispersión de valores reales contra predichos: si el modelo
fuera perfecto, todos los puntos caerían sobre la diagonal.

\begin{rcode}
grafico_real_pred <- ggplot(tiendas_prueba,
                            aes(x = pred_lm, y = ventas)) +
  geom_abline(slope = 1, intercept = 0, color = color_gris,
              linewidth = 0.8, linetype = "dashed") +
  geom_point(color = color_principal, size = 2, alpha = 0.8) +
  scale_x_continuous(labels = scales::dollar_format(suffix = "k")) +
  scale_y_continuous(labels = scales::dollar_format(suffix = "k")) +
  labs(title = "Ventas reales contra ventas predichas",
       subtitle = "80 tiendas de prueba; diagonal = predicción perfecta",
       x = "Ventas predichas (miles)", y = "Ventas reales (miles)") +
  theme_minimal() +
  theme(panel.grid.minor = element_blank())
grafico_real_pred
\end{rcode}

\begin{figure}[H]
\centering
\includegraphics[width=0.9\textwidth]{m07-real-vs-predicho}
\caption{Ventas reales contra predichas por la regresión lineal en las tiendas
de prueba. Los puntos se agrupan alrededor de la diagonal, sin patrones
sistemáticos.}
\label{fig:m07-real-vs-predicho}
\end{figure}

Busca dos cosas en esta gráfica: que los puntos estén cerca de la diagonal y
que el error no dependa del tamaño (por ejemplo, que el modelo no subestime
sistemáticamente las tiendas más grandes). Aquí ambas se cumplen.

\subsection{Usar el modelo: estimar las ventas de tiendas candidatas}

\begin{rcode}
candidatas <- tibble(
  tienda     = c("Local Plaza Norte", "Local Centro Histórico",
                 "Local Sur Express"),
  zona       = c("Norte", "Centro", "Sur"),
  tamano_m2  = c(900, 400, 250),
  publicidad = c(80, 50, 30),
  promocion  = c(1, 0, 1),
  trafico    = c(9000, 6500, 4000)
)

predict(modelo_lm, newdata = candidatas, interval = "prediction") %>%
  as_tibble() %>%
  mutate(tienda = candidatas$tienda, .before = 1) %>%
  mutate(across(where(is.numeric), ~ round(.x)))
\end{rcode}

\begin{rsalida}
# A tibble: 3 x 4
  tienda                   fit   lwr   upr
  <chr>                  <dbl> <dbl> <dbl>
1 Local Plaza Norte       1058   958  1159
2 Local Centro Histórico   630   529   730
3 Local Sur Express        456   354   557
\end{rsalida}

\codigo{interval = "prediction"} agrega un \termino{intervalo de predicción}
del 95\%: \codigo{fit} es la estimación puntual, y \codigo{lwr} y \codigo{upr}
los límites inferior y superior. Para
el local de Plaza Norte, el modelo estima ventas de \$1{,}058 mil al mes, y con
95\% de confianza entre \$958 mil y \$1{,}159 mil. Ese rango es lo que debes
llevar a la junta: si la renta solo es rentable con ventas mayores a \$1{,}100
mil, el local es una apuesta arriesgada aunque la estimación puntual se vea
bien. El local Sur Express, con ventas esperadas de \$456 mil, probablemente
no justifica la inversión.

\subsection{Comparación con random forest}

\begin{concepto}{Árboles de decisión y random forest}
Un \termino{árbol de decisión} hace preguntas sucesivas del tipo ``¿tamaño
mayor a 700 m²?'' y en cada hoja predice el promedio de las tiendas que llegan
ahí. Un \termino{random forest} (bosque aleatorio) entrena cientos de árboles,
cada uno con una muestra distinta de tiendas y de variables, y promedia sus
predicciones. Captura relaciones no lineales e interacciones sin que tengas
que especificarlas, pero es una ``caja negra'' más difícil de explicar.
\end{concepto}

\begin{rcode}
set.seed(123)
modelo_rf_ventas <- randomForest(
  ventas ~ publicidad + tamano_m2 + trafico + promocion + zona,
  data = tiendas_ent %>% mutate(zona = factor(zona)),
  ntree = 300,          # número de árboles
  importance = TRUE     # calcular importancia de variables
)
modelo_rf_ventas
\end{rcode}

\begin{rsalida}

Call:
...
               Type of random forest: regression
                     Number of trees: 300
No. of variables tried at each split: 1

          Mean of squared residuals: 7887.08
                    % Var explained: 80.2
\end{rsalida}

El bosque reporta un error \termino{OOB} (\emph{out-of-bag}): como cada árbol
se entrena con una muestra, las tiendas que quedaron fuera de esa muestra
sirven para evaluarlo, una especie de validación cruzada gratuita. Comparemos
ambos modelos en el conjunto de prueba:

\begin{rcode}
tiendas_prueba <- tiendas_prueba %>%
  mutate(pred_rf = predict(modelo_rf_ventas,
                           newdata = tiendas_prueba %>%
                             mutate(zona = factor(zona))))

bind_rows(
  metricas_regresion(tiendas_prueba$ventas, tiendas_prueba$pred_lm) %>%
    mutate(modelo = "Regresión lineal", .before = 1),
  metricas_regresion(tiendas_prueba$ventas, tiendas_prueba$pred_rf) %>%
    mutate(modelo = "Random forest", .before = 1)
) %>%
  mutate(across(where(is.numeric), ~ round(.x, 3)))
\end{rcode}

\begin{rsalida}
# A tibble: 2 x 5
  modelo             MAE  RMSE  MAPE    R2
  <chr>            <dbl> <dbl> <dbl> <dbl>
1 Regresión lineal  37.2  45.1  5.04 0.947
2 Random forest     67.7  85.3 10.5  0.809
\end{rsalida}

La regresión lineal gana con claridad: MAE de 37.2 contra 67.7 del random
forest, y R² de 0.947 contra 0.809. El bosque explica el 80.2\% de la varianza
según su estimación OOB, muy parecido a lo que obtuvo en la prueba (0.809): el
OOB es un buen adelanto del desempeño real.

\begin{negocio}{El modelo más sofisticado no siempre gana}
En esta simulación la relación es lineal por construcción, así que la
regresión lineal es el modelo ``correcto'' y le gana al bosque aleatorio. En
datos reales con umbrales y combinaciones (por ejemplo, la promoción solo
funciona en tiendas grandes), el bosque suele ganar. Por eso la práctica
profesional es: \textbf{probar un modelo simple y uno flexible, compararlos en
datos de prueba y quedarte con el más sencillo que cumpla}. Además, la
regresión lineal te da coeficientes que la dirección entiende.
\end{negocio}

\begin{hazlotu}
Quita \codigo{trafico} del modelo lineal y vuelve a evaluarlo. ¿Cuánto empeora
el MAE? ¿Cómo cambia el coeficiente de \codigo{tamano\_m2}? (Pista: al no estar
el tráfico, el tamaño absorbe parte de su efecto, porque ambas variables están
relacionadas.)
\end{hazlotu}

% ============================================================================
\section{Caso 3: clasificación para predecir el abandono de clientes}
\label{sec:m7-churn}

\begin{ejemplo}{¿A quién llama primero el equipo de retención?}
Una empresa de suscripción (piensa en un club de compras con membresía
mensual, un servicio de internet o un gimnasio) pierde cada mes a una parte de
sus clientes. El equipo de retención puede llamar a 150 clientes por semana y
ofrecerles un descuento. La pregunta: \emph{¿qué clientes tienen mayor
probabilidad de irse y, por lo tanto, a quiénes conviene llamar primero?} Es un
problema de \textbf{clasificación}: la respuesta es ``se va'' o ``se queda''.
\end{ejemplo}

El \termino{churn} (abandono, deserción o cancelación) es la proporción de
clientes que dejan de serlo en un periodo. Retener a un cliente suele costar
mucho menos que conseguir uno nuevo, por eso los modelos de churn son de los
más rentables en BI.

\subsection{Datos simulados de clientes}

Simulamos 2{,}000 clientes con variables típicas de un sistema de
suscripciones. El abandono se genera con reglas de negocio realistas: aumenta
con las quejas y los pagos tardíos, es mayor en los primeros meses, y disminuye
con la antigüedad, el uso del servicio y los planes más completos. El
\textbf{género y la edad no influyen} (los incluimos para una discusión de
ética más adelante).

\begin{rcode}
set.seed(123)
n_clientes <- 2000
clientes_churn <- tibble(
  cliente_id       = sprintf("C%04d", 1:n_clientes),
  antiguedad_meses = sample(1:72, n_clientes, replace = TRUE),
  plan             = sample(c("Básico", "Estándar", "Premium"),
                            n_clientes, replace = TRUE,
                            prob = c(0.5, 0.3, 0.2)),
  uso_mensual      = rpois(n_clientes, 8),    # visitas/usos al mes
  quejas_6m        = rpois(n_clientes, 0.8),  # quejas últimos 6 meses
  pagos_tardios    = rpois(n_clientes, 0.4),  # pagos atrasados en el año
  genero           = sample(c("F", "M"), n_clientes, replace = TRUE),
  edad             = sample(18:70, n_clientes, replace = TRUE)
) %>%
  mutate(
    # "Verdad" oculta: puntaje de riesgo (el modelo no la verá)
    efecto_plan = case_when(plan == "Básico"   ~ 0,
                            plan == "Estándar" ~ -0.6,
                            TRUE               ~ -1.3),
    riesgo = 2.0 - 0.05 * antiguedad_meses +
      1.0 * (antiguedad_meses <= 6) + 1.0 * quejas_6m -
      0.35 * uso_mensual + 0.8 * pagos_tardios + efecto_plan,
    prob_real = plogis(riesgo),               # convierte a probabilidad
    churn = if_else(runif(n_clientes) < prob_real, "Si", "No"),
    churn = factor(churn, levels = c("No", "Si")),
    plan  = factor(plan, levels = c("Básico", "Estándar", "Premium")),
    cuota_mensual = case_when(plan == "Básico"   ~ 199,
                              plan == "Estándar" ~ 349,
                              TRUE               ~ 599)
  ) %>%
  select(-efecto_plan, -riesgo, -prob_real)

glimpse(clientes_churn)
clientes_churn %>% count(churn) %>% mutate(pct = round(100 * n / sum(n), 1))
\end{rcode}

\begin{rsalida}
Rows: 2,000
Columns: 10
$ cliente_id       <chr> "C0001", "C0002", "C0003", "C0004", "C0005"~
$ antiguedad_meses <int> 31, 51, 14, 67, 42, 50, 43, 14, 25, 69, 57,~
$ plan             <fct> Estándar, Premium, Premium, Estándar, Básic~
$ uso_mensual      <int> 9, 8, 5, 14, 11, 14, 8, 10, 10, 8, 8, 9, 8,~
$ quejas_6m        <int> 1, 0, 0, 2, 1, 1, 1, 0, 0, 1, 0, 0, 3, 0, 0~
$ pagos_tardios    <int> 0, 3, 0, 1, 0, 0, 1, 0, 0, 0, 1, 1, 0, 0, 1~
$ genero           <chr> "M", "F", "M", "M", "F", "M", "M", "F", "M"~
$ edad             <int> 58, 44, 27, 57, 19, 43, 48, 36, 67, 18, 40,~
$ churn            <fct> No, No, No, No, No, No, No, No, No, No, No,~
$ cuota_mensual    <dbl> 349, 599, 599, 349, 199, 199, 349, 199, 349~
# A tibble: 2 x 3
  churn     n   pct
  <fct> <int> <dbl>
1 No     1545  77.2
2 Si      455  22.8
\end{rsalida}

\codigo{glimpse()} muestra las 10 columnas: 7 predictoras, el identificador,
la cuota mensual (que usaremos para calcular el valor del cliente) y la
variable objetivo \codigo{churn}. De los 2{,}000 clientes, 455 (22.8\%) se
fueron. Es una clase \termino{desbalanceada}: los que se quedan son más de
tres veces los que se van. Tenlo presente al leer la exactitud.

\begin{consejo}
Convierte la variable objetivo a \codigo{factor} con los niveles en un orden
explícito: \codigo{levels = c("No", "Si")}. Así R sabe que ``Si'' es la clase
de interés (la segunda), y \codigo{glm()} modelará la probabilidad de ``Si''.
Evitamos la tilde en ``Si'' para no tener problemas con nombres de columnas que
generan algunas funciones.
\end{consejo}

Antes de modelar, una mirada exploratoria rápida (como en el
Módulo~\ref{cap:eda}): ¿cómo cambia la tasa de abandono con las quejas?

\begin{rcode}
clientes_churn %>%
  mutate(quejas = if_else(quejas_6m >= 3, "3 o más",
                          as.character(quejas_6m))) %>%
  group_by(quejas) %>%
  summarise(clientes = n(),
            tasa_churn = round(100 * mean(churn == "Si"), 1),
            .groups = "drop")
\end{rcode}

\begin{rsalida}
# A tibble: 4 x 3
  quejas  clientes tasa_churn
  <chr>      <int>      <dbl>
1 0            903       13.4
2 1            747       22.6
3 2            270       42.6
4 3 o más       80       62.5
\end{rsalida}

La relación es clarísima: sin quejas, solo 13.4\% abandona; con 3 o más
quejas, 62.5\%. Este tipo de tabla es la mejor forma de verificar, antes de
modelar, que hay señal en los datos.

Dividimos en entrenamiento (70\%) y prueba (30\%) de forma
\textbf{estratificada} con la función \codigo{createDataPartition()}, para que ambos
conjuntos tengan la misma proporción de abandono:

\begin{rcode}
set.seed(123)
filas_ent <- createDataPartition(clientes_churn$churn, p = 0.7,
                                 list = FALSE)
churn_ent    <- clientes_churn[filas_ent, ]
churn_prueba <- clientes_churn[-filas_ent, ]

# Misma proporción de "Si" en ambos conjuntos
round(prop.table(table(churn_ent$churn)), 3)
round(prop.table(table(churn_prueba$churn)), 3)
\end{rcode}

\begin{rsalida}

   No    Si 
0.772 0.228 

   No    Si 
0.773 0.227 
\end{rsalida}

\subsection{Regresión logística}

\begin{concepto}{Regresión logística}
La \termino{regresión logística} es la ``regresión lineal de la
clasificación''. En lugar de predecir un número cualquiera, predice una
\textbf{probabilidad} entre 0 y 1 (por ejemplo, 0.73 = 73\% de probabilidad de
abandono). En R se ajusta con \codigo{glm(..., family = binomial)}. Es rápida,
estable y, sobre todo, \textbf{explicable}: cada coeficiente dice cómo cambia
el riesgo.
\end{concepto}

\begin{rcode}
formula_churn <- churn ~ antiguedad_meses + plan + uso_mensual +
  quejas_6m + pagos_tardios + genero + edad

modelo_log <- glm(formula_churn, data = churn_ent, family = binomial)
summary(modelo_log)
\end{rcode}

\begin{rsalida}

Call:
glm(formula = formula_churn, family = binomial, data = churn_ent)

Coefficients:
                  Estimate Std. Error z value Pr(>|z|)    
(Intercept)       2.114728   0.385699   5.483 4.19e-08 ***
antiguedad_meses -0.054049   0.004438 -12.178  < 2e-16 ***
planEstándar     -0.448662   0.175005  -2.564   0.0104 *  
planPremium      -1.247963   0.234692  -5.317 1.05e-07 ***
uso_mensual      -0.310629   0.032282  -9.622  < 2e-16 ***
quejas_6m         1.048309   0.091879  11.410  < 2e-16 ***
pagos_tardios     0.528275   0.121604   4.344 1.40e-05 ***
generoM           0.019595   0.154767   0.127   0.8992    
edad             -0.003247   0.005179  -0.627   0.5306    
---
Signif. codes:  0 '***' 0.001 '**' 0.01 '*' 0.05 '.' 0.1 ' ' 1

(Dispersion parameter for binomial family taken to be 1)

    Null deviance: 1503.2  on 1400  degrees of freedom
Residual deviance: 1063.0  on 1392  degrees of freedom
AIC: 1081

Number of Fisher Scoring iterations: 5
\end{rsalida}

Los coeficientes de la logística están en una escala poco intuitiva
(\emph{log-odds}). Para explicarlos al negocio, conviértelos en
\termino{razones de momios} (\emph{odds ratios}) con \codigo{exp()}:

\begin{rcode}
tibble(variable   = names(coef(modelo_log)),
       odds_ratio = round(exp(coef(modelo_log)), 3)) %>%
  mutate(efecto = case_when(
    odds_ratio > 1 ~ paste0("+", round(100 * (odds_ratio - 1)), "% riesgo"),
    TRUE           ~ paste0(round(100 * (odds_ratio - 1)), "% riesgo")
  )) %>%
  filter(variable != "(Intercept)")
\end{rcode}

\begin{rsalida}
# A tibble: 8 x 3
  variable         odds_ratio efecto      
  <chr>                 <dbl> <chr>       
1 antiguedad_meses      0.947 -5% riesgo  
2 planEstándar          0.638 -36% riesgo 
3 planPremium           0.287 -71% riesgo 
4 uso_mensual           0.733 -27% riesgo 
5 quejas_6m             2.85  +185% riesgo
6 pagos_tardios         1.70  +70% riesgo 
7 generoM               1.02  +2% riesgo  
8 edad                  0.997 0% riesgo   
\end{rsalida}

\begin{concepto}{Cómo leer una razón de momios}
Los \emph{momios} (\emph{odds}) son la probabilidad de que algo pase dividida
entre la probabilidad de que no pase. Una razón de momios de 1 significa ``no
cambia el riesgo''; mayor a 1, lo aumenta; menor a 1, lo reduce. Se interpreta
\emph{por cada unidad} de la variable y manteniendo las demás constantes.
\end{concepto}

Traducido al negocio:
\begin{itemize}
  \item \textbf{Quejas:} cada queja adicional multiplica los momios de
        abandono por 2.85 (casi los triplica). Es la palanca más fuerte:
        resolver bien una queja vale mucho.
  \item \textbf{Pagos tardíos:} cada pago atrasado aumenta los momios 70\%.
        Es una señal temprana que cobranza ya tiene y puede compartir.
  \item \textbf{Uso:} cada uso adicional al mes reduce los momios 27\%. Los
        clientes que usan el servicio se quedan.
  \item \textbf{Antigüedad:} cada mes adicional reduce los momios 5\%; en un
        año, eso se acumula ($0.947^{12} \approx 0.52$, casi la mitad).
  \item \textbf{Plan:} frente al Básico, el Premium reduce los momios 71\% y
        el Estándar 36\%.
  \item \textbf{Género y edad:} razones cercanas a 1 y valores p altos (0.90
        y 0.53 en el resumen): no influyen, tal como los simulamos.
\end{itemize}
Observa también que el modelo no ``descubrió'' el riesgo extra de los primeros
6 meses, porque no le dimos una variable que lo represente; la antigüedad
lineal lo absorbe en parte. Crear variables como
\codigo{cliente\_nuevo = antiguedad\_meses <= 6} es el tipo de mejora que se
logra entendiendo el negocio (ejercicio~\ref{ej:m7-6}).

\subsection{Predicciones y matriz de confusión}

El modelo entrega probabilidades. Para convertirlas en una decisión (llamar o
no llamar) necesitamos un \termino{umbral}: si la probabilidad es mayor o igual
al umbral, clasificamos como ``Si''. Empecemos con el umbral por defecto,
0.5:

\begin{rcode}
churn_prueba <- churn_prueba %>%
  mutate(prob_log = predict(modelo_log, newdata = churn_prueba,
                            type = "response"),   # probabilidad de "Si"
         pred_log = factor(if_else(prob_log >= 0.5, "Si", "No"),
                           levels = c("No", "Si")))

churn_prueba %>%
  select(cliente_id, antiguedad_meses, quejas_6m, prob_log, pred_log,
         churn) %>%
  mutate(prob_log = round(prob_log, 3)) %>%
  head(5)

# Matriz de confusión: filas = predicción, columnas = realidad
matriz <- table(Prediccion = churn_prueba$pred_log,
                Real = churn_prueba$churn)
matriz
\end{rcode}

\begin{rsalida}
# A tibble: 5 x 6
  cliente_id antiguedad_meses quejas_6m prob_log pred_log churn
  <chr>                 <int>     <int>    <dbl> <fct>    <fct>
1 C0003                    14         0    0.181 No       No   
2 C0007                    43         1    0.154 No       No   
3 C0012                     9         0    0.119 No       No   
4 C0014                    26         0    0.224 No       No   
5 C0015                     7         0    0.416 No       Si   
          Real
Prediccion  No  Si
        No 434  69
        Si  29  67
\end{rsalida}

\begin{concepto}{Matriz de confusión}
Compara lo que el modelo predijo con lo que realmente pasó:

\begin{center}
\small
\begin{tabular}{lcc}
\toprule
 & \textbf{Real: No se fue} & \textbf{Real: Sí se fue} \\
\midrule
\textbf{Predicción: No} & Verdadero negativo (VN) & Falso negativo (FN) \\
\textbf{Predicción: Sí} & Falso positivo (FP) & Verdadero positivo (VP) \\
\bottomrule
\end{tabular}
\end{center}

\begin{itemize}
  \item \textbf{Falso positivo:} llamas a un cliente que no se iba a ir
        (desperdicias un descuento).
  \item \textbf{Falso negativo:} no llamas a un cliente que sí se va (lo
        pierdes). Suele ser el error más caro.
\end{itemize}
\end{concepto}

En nuestros 599 clientes de prueba: el modelo acertó 434 clientes que se
quedaron (VN) y 67 que se fueron (VP). Se equivocó con 29 clientes que marcó
como ``se va'' y se quedaron (FP) y, lo más grave, dejó pasar a 69 clientes que
se fueron sin que los detectara (FN). Mira también el cliente C0015 en la
tabla de arriba: tenía 41.6\% de probabilidad y sí se fue; con el umbral de 0.5
no lo habríamos llamado.

\subsection{Métricas de clasificación}
\label{sec:m7-metricas-clas}

De la matriz de confusión salen todas las métricas. Las calculamos a mano para
entender qué significa cada una:

\begin{rcode}
metricas_clasificacion <- function(prediccion, real, positivo = "Si") {
  vp <- sum(prediccion == positivo & real == positivo)
  vn <- sum(prediccion != positivo & real != positivo)
  fp <- sum(prediccion == positivo & real != positivo)
  fn <- sum(prediccion != positivo & real == positivo)
  precision    <- vp / (vp + fp)
  sensibilidad <- vp / (vp + fn)
  tibble(
    exactitud    = (vp + vn) / (vp + vn + fp + fn),
    precision    = precision,
    sensibilidad = sensibilidad,
    especificidad = vn / (vn + fp),
    F1           = 2 * precision * sensibilidad /
                   (precision + sensibilidad)
  )
}

metricas_clasificacion(churn_prueba$pred_log, churn_prueba$churn) %>%
  mutate(across(everything(), ~ round(.x, 3)))
\end{rcode}

\begin{rsalida}
# A tibble: 1 x 5
  exactitud precision sensibilidad especificidad    F1
      <dbl>     <dbl>        <dbl>         <dbl> <dbl>
1     0.836     0.698        0.493         0.937 0.578
\end{rsalida}

\begin{table}[H]
\centering
\small
\caption{Métricas de clasificación y la pregunta de negocio que responden.}
\label{tab:m7-metricas-clas}
\begin{tabularx}{\textwidth}{>{\raggedright\arraybackslash}p{2.6cm}
                              >{\raggedright\arraybackslash}p{3.3cm}
                              >{\raggedright\arraybackslash}X}
\toprule
\textbf{Métrica} & \textbf{Fórmula} & \textbf{Pregunta que responde} \\
\midrule
Exactitud (\emph{accuracy}) & $(VP + VN) / \text{total}$ &
  ¿Qué proporción de clientes clasificamos bien? Engañosa si una clase es
  rara. \\
Precisión (\emph{precision}) & $VP / (VP + FP)$ &
  De los clientes que marcamos como ``se va'', ¿cuántos realmente se iban?
  (eficiencia de las llamadas). \\
Sensibilidad (\emph{recall}) & $VP / (VP + FN)$ &
  De los clientes que se fueron, ¿a cuántos detectamos? (cobertura). \\
Especificidad & $VN / (VN + FP)$ &
  De los que se quedaron, ¿a cuántos dejamos tranquilos? \\
F1 & $2 \cdot \dfrac{\text{prec.} \cdot \text{sens.}}{\text{prec.} + \text{sens.}}$ &
  Equilibrio entre precisión y sensibilidad en un solo número. \\
\bottomrule
\end{tabularx}
\end{table}

La exactitud es de 83.6\%, que suena bien, pero la sensibilidad es de solo
49.3\%: con el umbral de 0.5 detectamos a menos de la mitad de los clientes que
se van. La precisión de 69.8\% dice que de cada 10 llamadas, 7 serían a
clientes realmente en riesgo. Recuerda que predecir ``nadie se va'' ya daría
77.3\% de exactitud (la proporción de ``No'' en la prueba), así que la exactitud
por sí sola dice poco (sección~\ref{sec:m7-errores}).

\subsection{Curva ROC y AUC}

Cada umbral produce una matriz de confusión distinta. La \termino{curva ROC}
muestra todos los umbrales a la vez: en el eje horizontal, la tasa de falsos
positivos ($1 -$ especificidad); en el vertical, la sensibilidad. El
\termino{AUC} (área bajo la curva) resume la curva en un número: la
probabilidad de que, al tomar un cliente que se fue y uno que se quedó al azar,
el modelo le asigne mayor riesgo al que se fue. Un AUC de 0.5 es adivinar con
una moneda; 1 es perfecto; en churn, de 0.75 a 0.85 ya es un buen modelo.

Calculamos la curva a mano, probando 101 umbrales:

\begin{rcode}
curva_roc <- function(prob, real, positivo = "Si") {
  map_dfr(seq(0, 1, by = 0.01), function(u) {
    pred <- prob >= u
    tibble(umbral = u,
           sensibilidad = sum(pred & real == positivo) /
                          sum(real == positivo),
           tasa_fp = sum(pred & real != positivo) / sum(real != positivo))
  })
}

# AUC como probabilidad de ordenar bien un par (positivo, negativo)
calcular_auc <- function(prob, real, positivo = "Si") {
  p_pos <- prob[real == positivo]
  p_neg <- prob[real != positivo]
  mean(outer(p_pos, p_neg, ">") + 0.5 * outer(p_pos, p_neg, "=="))
}

roc_log <- curva_roc(churn_prueba$prob_log, churn_prueba$churn)
auc_log <- calcular_auc(churn_prueba$prob_log, churn_prueba$churn)
round(auc_log, 3)
# round() evita problemas de precisión al comparar decimales
roc_log %>% filter(round(umbral, 2) %in% c(0.1, 0.2, 0.3, 0.5, 0.7))
\end{rcode}

\begin{rsalida}
[1] 0.86
# A tibble: 5 x 3
  umbral sensibilidad tasa_fp
   <dbl>        <dbl>   <dbl>
1    0.1        0.941  0.460 
2    0.2        0.801  0.298 
3    0.3        0.691  0.173 
4    0.5        0.493  0.0626
5    0.7        0.331  0.0151
\end{rsalida}

El AUC es de 0.86: si tomas al azar un cliente que se fue y uno que se quedó,
en 86\% de los casos el modelo le asigna más riesgo al que se fue. La tabla
muestra el dilema del umbral: con 0.5 detectamos 49.3\% de los que se van
molestando solo a 6.3\% de los que se quedan; con 0.2 detectamos 80.1\%, pero
molestamos a 29.8\%; con 0.1, 94.1\% y 46\%. En el otro extremo, con 0.7 casi no molestamos a nadie (1.5\%), pero solo
detectamos a un tercio (33.1\%) de los que se van.

\begin{rcode}
grafico_roc <- ggplot(roc_log, aes(x = tasa_fp, y = sensibilidad)) +
  geom_abline(slope = 1, intercept = 0, color = color_gris,
              linetype = "dashed", linewidth = 0.8) +
  geom_path(color = color_principal, linewidth = 0.8) +
  geom_point(data = filter(roc_log, umbral %in% c(0.2, 0.5)),
             color = color_resalte, size = 3) +
  geom_text(data = filter(roc_log, umbral %in% c(0.2, 0.5)),
            aes(label = paste("umbral", umbral)), color = color_resalte,
            hjust = -0.2, vjust = 1.2, size = 3.5) +
  annotate("text", x = 0.75, y = 0.25,
           label = paste("AUC =", round(auc_log, 3)), size = 4.5) +
  scale_x_continuous(labels = scales::percent) +
  scale_y_continuous(labels = scales::percent) +
  coord_equal() +
  labs(title = "Curva ROC de la regresión logística",
       subtitle = "Clientes de prueba; la diagonal equivale a adivinar",
       x = "Tasa de falsos positivos (1 - especificidad)",
       y = "Sensibilidad") +
  theme_minimal() +
  theme(panel.grid.minor = element_blank())
grafico_roc
\end{rcode}

\begin{figure}[H]
\centering
\includegraphics[width=0.9\textwidth]{m07-roc}
\caption{Curva ROC del modelo de abandono. Bajar el umbral de 0.5 a 0.2 sube
mucho la sensibilidad a cambio de más falsos positivos.}
\label{fig:m07-roc}
\end{figure}

\subsection{Elegir el umbral según el costo de negocio}
\label{sec:m7-umbral}

El umbral de 0.5 no tiene nada de especial. La decisión correcta depende de
cuánto cuesta cada error. Supongamos estos valores (en pesos):

\begin{table}[H]
\centering
\small
\caption{Supuestos de costos para la campaña de retención.}
\label{tab:m7-costos}
\begin{tabularx}{\textwidth}{>{\raggedright\arraybackslash}p{4.2cm}
                              >{\raggedright\arraybackslash}X r}
\toprule
\textbf{Resultado} & \textbf{Qué pasa} & \textbf{Costo} \\
\midrule
Verdadero positivo (VP) & Llamamos y ofrecemos descuento; la mitad de ellos se
  queda y la otra mitad se va de todos modos. & $300 + 0.5 \times 3{,}000 = 1{,}800$ \\
Falso positivo (FP) & Damos un descuento a quien no se iba a ir. & 300 \\
Falso negativo (FN) & No llamamos y el cliente se va: perdemos su margen de
  un año. & 3{,}000 \\
Verdadero negativo (VN) & No hacemos nada y el cliente se queda. & 0 \\
\bottomrule
\end{tabularx}
\end{table}

Con esos costos, calculamos el costo total de la campaña para cada umbral y
buscamos el mínimo:

\begin{rcode}
costo_vp <- 300 + 0.5 * 3000   # descuento + clientes que igual se van
costo_fp <- 300                # descuento desperdiciado
costo_fn <- 3000               # cliente perdido sin intentar retenerlo

costos_umbral <- map_dfr(seq(0.05, 0.95, by = 0.05), function(u) {
  pred <- churn_prueba$prob_log >= u
  real <- churn_prueba$churn == "Si"
  tibble(umbral = u,
         llamadas = sum(pred),
         vp = sum(pred & real), fp = sum(pred & !real),
         fn = sum(!pred & real),
         costo_total = vp * costo_vp + fp * costo_fp + fn * costo_fn)
})

costos_umbral %>% arrange(costo_total) %>% head(3)
costos_umbral %>% filter(umbral %in% c(0.5, 0.95))
\end{rcode}

\begin{rsalida}
# A tibble: 3 x 6
  umbral llamadas    vp    fp    fn costo_total
   <dbl>    <int> <int> <int> <int>       <dbl>
1   0.15      294   120   174    16      316200
2   0.1       341   128   213     8      318300
3   0.2       247   109   138    27      318600
# A tibble: 2 x 6
  umbral llamadas    vp    fp    fn costo_total
   <dbl>    <int> <int> <int> <int>       <dbl>
1   0.5        96    67    29    69      336300
2   0.95        4     4     0   132      403200
\end{rsalida}

Con estos costos, el mejor umbral es 0.15: llamaríamos a 294 clientes de
prueba, detectaríamos a 120 de los 136 que se van y el costo total sería
\$316{,}200. Con el umbral por defecto de 0.5 el costo sube a \$336{,}300,
y casi sin llamar a nadie (umbral 0.95) a \$403{,}200. Elegir bien el umbral
ahorra \$20{,}100 en solo 599 clientes; escalado a una cartera de 100{,}000
clientes, son más de 3 millones de pesos.

\begin{rcode}
grafico_costos <- ggplot(costos_umbral, aes(umbral, costo_total)) +
  geom_line(color = color_principal, linewidth = 0.8) +
  geom_point(color = color_principal, size = 1.5) +
  geom_point(data = slice_min(costos_umbral, costo_total, n = 1),
             color = color_resalte, size = 4) +
  scale_y_continuous(labels = scales::dollar) +
  labs(title = "Costo total de la campaña según el umbral",
       subtitle = "Clientes de prueba; en naranja, el umbral óptimo",
       x = "Umbral de probabilidad para llamar al cliente",
       y = "Costo total") +
  theme_minimal() +
  theme(panel.grid.minor = element_blank())
grafico_costos
\end{rcode}

\begin{figure}[H]
\centering
\includegraphics[width=0.9\textwidth]{m07-costo-umbral}
\caption{Costo total de la campaña de retención para cada umbral. Con errores
de costos tan distintos, conviene un umbral mucho más bajo que 0.5.}
\label{fig:m07-costo-umbral}
\end{figure}

\begin{negocio}{El umbral es una decisión de negocio, no técnica}
Cuando el costo de perder un cliente es mucho mayor que el de un descuento, el
umbral óptimo baja: conviene llamar a más gente aunque muchos no se fueran a ir.
Si el equipo de retención tiene capacidad limitada (150 llamadas por semana),
el umbral lo fija la capacidad: llamas a los 150 de mayor probabilidad. Lo
importante es que el modelo \emph{ordena} bien a los clientes por riesgo, y
eso lo mide el AUC.
\end{negocio}

\begin{hazlotu}
Repite el cálculo de costos suponiendo que el descuento cuesta \$800 en lugar
de \$300. ¿Sube o baja el umbral óptimo? ¿Por qué?
\end{hazlotu}

\subsection{Árbol de decisión con \texorpdfstring{\codigo{rpart}}{rpart}}

Un \termino{árbol de decisión} produce reglas del tipo ``si antigüedad < 20
meses y uso < 8, entonces riesgo alto''. Su gran ventaja es que cualquier
gerente las entiende. Para que el árbol sea legible limitamos su complejidad
con el parámetro \codigo{cp} (\emph{complexity parameter}): un corte solo se
hace si mejora el ajuste al menos en esa proporción.

\begin{rcode}
set.seed(123)
modelo_arbol <- rpart(formula_churn, data = churn_ent, method = "class",
                      control = rpart.control(cp = 0.02))
modelo_arbol
\end{rcode}

\begin{rsalida}
n= 1401 

node), split, n, loss, yval, (yprob)
      * denotes terminal node

 1) root 1401 319 No (0.7723055 0.2276945)  
   2) antiguedad_meses>=20.5 1014 151 No (0.8510848 0.1489152) *
   3) antiguedad_meses< 20.5 387 168 No (0.5658915 0.4341085)  
     6) uso_mensual>=7.5 211  63 No (0.7014218 0.2985782) *
     7) uso_mensual< 7.5 176  71 Si (0.4034091 0.5965909)  
      14) quejas_6m< 1.5 134  64 Si (0.4776119 0.5223881)  
        28) plan=Estándar,Premium 66  25 No (0.6212121 0.3787879) *
        29) plan=Básico 68  23 Si (0.3382353 0.6617647) *
      15) quejas_6m>=1.5 42   7 Si (0.1666667 0.8333333) *
\end{rsalida}

Cada línea es un nodo: número, regla, clientes que llegan, clientes mal
clasificados, clase predicha y proporciones (No, Si). Las hojas se marcan con
\codigo{*}. Leídas como reglas de negocio:

\begin{enumerate}
  \item \textbf{Antigüedad $\geq$ 21 meses} (1{,}014 clientes): riesgo bajo,
        14.9\% abandona.
  \item \textbf{Antigüedad < 21 meses y uso $\geq$ 8} (211): riesgo medio,
        29.9\%.
  \item \textbf{Antigüedad < 21, uso < 8 y 2 o más quejas} (42): riesgo muy
        alto, 83.3\%.
  \item \textbf{Antigüedad < 21, uso < 8, menos de 2 quejas y plan Básico}
        (68): riesgo alto, 66.2\%.
  \item Lo mismo pero con plan Estándar o Premium (66): 37.9\%.
\end{enumerate}

Estas reglas se pueden pegar tal cual en una presentación o programar en SQL.
Coinciden con lo que dijo la logística: antigüedad, uso, quejas y plan.

Para dibujar el árbol, el paquete \paquete{rpart.plot} produce diagramas muy
claros. No viene con R; si lo instalas, bastan dos líneas:

\begin{rnoejecutar}
install.packages("rpart.plot")
library(rpart.plot)
rpart.plot(modelo_arbol, type = 4, extra = 104,
           box.palette = "BuOr", main = "Árbol de abandono")
\end{rnoejecutar}

Sin paquetes extra, \codigo{plot(modelo\_arbol); text(modelo\_arbol)} dibuja
una versión sencilla con gráficos base. ¿Qué variables usa más el árbol?

\begin{rcode}
round(modelo_arbol$variable.importance, 1)

churn_prueba <- churn_prueba %>%
  mutate(prob_arbol = predict(modelo_arbol, newdata = churn_prueba,
                              type = "prob")[, "Si"])
round(calcular_auc(churn_prueba$prob_arbol, churn_prueba$churn), 3)
\end{rcode}

\begin{rsalida}
antiguedad_meses      uso_mensual        quejas_6m             plan 
            46.1             17.4              7.0              5.4 
            edad           genero    pagos_tardios 
             0.6              0.4              0.2 
[1] 0.674
\end{rsalida}

La antigüedad es, por mucho, la variable más importante para el árbol,
seguida por el uso. Pero su AUC en prueba es de solo 0.674, bastante peor que
el 0.86 de la logística: con apenas 5 hojas, el árbol asigna la misma
probabilidad a cientos de clientes y no puede ordenarlos con detalle. Es el
precio de la simplicidad.

\subsection{Random forest}

\begin{rcode}
set.seed(123)
modelo_rf <- randomForest(formula_churn, data = churn_ent,
                          ntree = 300, importance = TRUE)
modelo_rf
\end{rcode}

\begin{rsalida}

Call:
...
               Type of random forest: classification
                     Number of trees: 300
No. of variables tried at each split: 2

        OOB estimate of  error rate: 19.49%
Confusion matrix:
     No  Si class.error
No 1008  74  0.06839187
Si  199 120  0.62382445
\end{rsalida}

El error OOB es de 19.5\%, pero la matriz de confusión OOB revela el mismo
problema de antes: el bosque clasifica bien a los que se quedan (6.8\% de error)
y mal a los que se van (62.4\% de error), porque usa el umbral de 0.5 por
defecto. Como vimos, para churn conviene trabajar con las probabilidades y
elegir el umbral.

\begin{rcode}
importancia_rf <- importance(modelo_rf) %>%
  as.data.frame() %>%
  rownames_to_column("variable") %>%
  as_tibble() %>%
  select(variable, MeanDecreaseAccuracy) %>%
  arrange(desc(MeanDecreaseAccuracy))
importancia_rf %>% mutate(MeanDecreaseAccuracy =
                            round(MeanDecreaseAccuracy, 1))

grafico_importancia <- ggplot(importancia_rf,
       aes(x = MeanDecreaseAccuracy,
           y = reorder(variable, MeanDecreaseAccuracy))) +
  geom_col(fill = color_principal, width = 0.7) +
  labs(title = "¿Qué variables usa el bosque para predecir el abandono?",
       subtitle = "Pérdida de exactitud al desordenar cada variable",
       x = "Importancia (disminución media de exactitud)", y = NULL) +
  theme_minimal() +
  theme(panel.grid.minor = element_blank(),
        panel.grid.major.y = element_blank())
grafico_importancia
\end{rcode}

\begin{rsalida}
# A tibble: 7 x 2
  variable         MeanDecreaseAccuracy
  <chr>                           <dbl>
1 antiguedad_meses                 32.6
2 quejas_6m                        31.3
3 uso_mensual                      23  
4 plan                              6.9
5 pagos_tardios                     4.7
6 genero                            3  
7 edad                              2.4
\end{rsalida}

\begin{figure}[H]
\centering
\includegraphics[width=0.9\textwidth]{m07-importancia}
\caption{Importancia de variables del random forest. Antigüedad, quejas y uso
dominan; género y edad no aportan.}
\label{fig:m07-importancia}
\end{figure}

\begin{concepto}{Importancia por permutación}
\codigo{MeanDecreaseAccuracy} responde: ``si desordeno al azar los valores de
esta variable (rompiendo su relación con el abandono), ¿cuánto empeora el
modelo?''. Si casi no empeora, la variable no aportaba. Es una forma de abrir
un poco la caja negra.
\end{concepto}

El bosque coincide con los demás modelos: antigüedad (32.6), quejas (31.3) y
uso (23.0) son las variables clave; plan y pagos tardíos aportan algo; género y
edad aportan muy poco. Que distintos algoritmos coincidan en las variables
importantes da confianza en el hallazgo.

\subsection{Comparar modelos con \texorpdfstring{\paquete{caret}}{caret}}

Hasta ahora cada algoritmo tenía su propia sintaxis. El paquete
\paquete{caret} ofrece una interfaz común, \codigo{train()}, para más de 200
algoritmos, y se encarga de la validación cruzada y del ajuste de parámetros.
Comparemos los tres modelos con validación cruzada de 5 pliegues usando el AUC
(que \paquete{caret} llama \codigo{ROC}) como criterio.

\begin{rcode}
# caret toma como clase positiva el PRIMER nivel del factor
churn_ent_caret <- churn_ent %>%
  mutate(churn = factor(churn, levels = c("Si", "No")))

control_cv <- trainControl(
  method = "cv", number = 5,           # validación cruzada de 5 pliegues
  classProbs = TRUE,                   # calcular probabilidades
  summaryFunction = twoClassSummary    # reportar ROC, Sens y Spec
)

set.seed(123)
cv_log <- train(formula_churn, data = churn_ent_caret, method = "glm",
                family = binomial, trControl = control_cv, metric = "ROC")
set.seed(123)
cv_arbol <- train(formula_churn, data = churn_ent_caret,
                  method = "rpart", trControl = control_cv,
                  metric = "ROC", tuneLength = 5)   # prueba 5 valores de cp
set.seed(123)
cv_rf <- train(formula_churn, data = churn_ent_caret, method = "rf",
               trControl = control_cv, metric = "ROC", ntree = 200,
               tuneGrid = data.frame(mtry = 2:4))   # variables por corte

cv_rf
\end{rcode}

\begin{rsalida}
Random Forest 

1401 samples
   7 predictor
   2 classes: 'Si', 'No' 

No pre-processing
Resampling: Cross-Validated (5 fold) 
Summary of sample sizes: 1121, 1121, 1120, 1122, 1120 
Resampling results across tuning parameters:

  mtry  ROC        Sens       Spec     
  2     0.8187057  0.3446925  0.9454600
  3     0.8115210  0.4198413  0.9149727
  4     0.8048380  0.4293155  0.9094342

ROC was used to select the optimal model using the largest value.
The final value used for the model was mtry = 2.
\end{rsalida}

\codigo{train()} probó tres valores de \codigo{mtry} (cuántas variables
considera el bosque en cada corte) y eligió el de mayor AUC promedio en los 5
pliegues. Ahora juntamos los tres modelos en una tabla:

\begin{rcode}
comparacion <- resamples(list(Logistica = cv_log, Arbol = cv_arbol,
                              RandomForest = cv_rf))

summary(comparacion)$statistics$ROC[, c("Min.", "Mean", "Max.")] %>%
  round(3)
\end{rcode}

\begin{rsalida}
              Min.  Mean  Max.
Logistica    0.828 0.843 0.882
Arbol        0.657 0.691 0.735
RandomForest 0.795 0.819 0.842
\end{rsalida}

Y confirmamos con los datos de prueba, que ningún modelo ha visto:

\begin{rcode}
churn_prueba <- churn_prueba %>%
  mutate(prob_rf = predict(modelo_rf, newdata = churn_prueba,
                           type = "prob")[, "Si"])

tibble(
  modelo = c("Regresión logística", "Árbol de decisión",
             "Random forest"),
  auc_prueba = c(calcular_auc(churn_prueba$prob_log, churn_prueba$churn),
                 calcular_auc(churn_prueba$prob_arbol, churn_prueba$churn),
                 calcular_auc(churn_prueba$prob_rf, churn_prueba$churn))
) %>%
  mutate(auc_prueba = round(auc_prueba, 3)) %>%
  arrange(desc(auc_prueba))
\end{rcode}

\begin{rsalida}
# A tibble: 3 x 2
  modelo              auc_prueba
  <chr>                    <dbl>
1 Regresión logística      0.86 
2 Random forest            0.846
3 Árbol de decisión        0.674
\end{rsalida}

Validación cruzada y prueba cuentan la misma historia: la regresión logística
gana (AUC de 0.843 en validación cruzada y 0.86 en prueba), seguida muy de
cerca por el random forest (0.819 y 0.846), mientras que el árbol sencillo
queda muy atrás (0.691 y 0.674). Como en la regresión de ventas, la relación
simulada es casi lineal y el modelo simple gana. Con la logística tenemos,
además, odds ratios que el negocio entiende: es el modelo que llevaríamos a
producción.

\subsection{Del modelo a la acción: lista priorizada de clientes en riesgo}

El entregable final para el equipo de retención no es un AUC: es una
\textbf{lista de clientes ordenada por prioridad}. Una buena práctica es
ordenar por \termino{pérdida esperada}: probabilidad de abandono $\times$
valor del cliente (aquí, su cuota de un año). Así, un cliente Premium con 50\%
de riesgo puede ir antes que uno Básico con 70\%.

\begin{rcode}
lista_retencion <- churn_prueba %>%
  mutate(valor_anual     = cuota_mensual * 12,
         perdida_esperada = prob_log * valor_anual) %>%
  arrange(desc(perdida_esperada)) %>%
  mutate(prioridad = row_number()) %>%
  select(prioridad, cliente_id, plan, quejas_6m, prob_log,
         perdida_esperada)

lista_retencion %>%
  mutate(prob_log = round(prob_log, 2),
         perdida_esperada = round(perdida_esperada)) %>%
  head(8)
\end{rcode}

\begin{rsalida}
# A tibble: 8 x 6
  prioridad cliente_id plan     quejas_6m prob_log perdida_esperada
      <int> <chr>      <fct>        <int>    <dbl>            <dbl>
1         1 C0259      Premium          2     0.88             6313
2         2 C1336      Premium          2     0.87             6261
3         3 C0609      Premium          1     0.67             4838
4         4 C1058      Premium          1     0.66             4768
5         5 C1912      Premium          2     0.57             4085
6         6 C1072      Estándar         3     0.97             4049
7         7 C1354      Premium          2     0.56             4029
8         8 C0795      Premium          2     0.55             3973
\end{rsalida}

Los primeros lugares los ocupan clientes Premium con probabilidad alta: su
cuota de \$599 hace que perderlos cueste más. El cliente C1072 tiene la mayor
probabilidad de la tabla (0.97), pero como es Estándar queda en el sexto lugar.
Esta lista, exportada a CSV cada semana, es lo que realmente usa el call
center (sección~\ref{sec:m7-produccion}).

\begin{hazlotu}
Agrega a la lista una columna \codigo{motivo\_principal} con
\codigo{case\_when()}: ``Quejas'' si \codigo{quejas\_6m >= 2}, ``Cliente
nuevo'' si \codigo{antiguedad\_meses <= 6}, ``Poco uso'' si
\codigo{uso\_mensual < 5} y ``Otro'' en los demás casos. El equipo de retención
podrá adaptar su guion según el motivo.
\end{hazlotu}

% ============================================================================
\section{Caso 4: pronóstico de demanda con series de tiempo}
\label{sec:m7-series}

\begin{ejemplo}{¿Cuántas unidades pedimos para diciembre?}
El área de compras de la categoría de electrónica debe colocar sus pedidos con
dos meses de anticipación. Si pide de menos en noviembre y diciembre (Buen Fin
y Navidad) pierde ventas; si pide de más, se queda con inventario que en enero
nadie quiere. La pregunta: \emph{¿cuántas unidades venderemos cada mes del
próximo año, y con qué margen de error?}
\end{ejemplo}

Una \termino{serie de tiempo} es una secuencia de mediciones ordenadas en el
tiempo a intervalos regulares (ventas diarias, demanda mensual). Sus
componentes típicos son:

\begin{itemize}
  \item \termino{Tendencia}: la dirección de largo plazo (el negocio crece).
  \item \termino{Estacionalidad}: patrones que se repiten cada año (picos en
        noviembre y diciembre, caída en enero y febrero).
  \item \termino{Ruido} (componente irregular): lo que queda, variaciones que
        no siguen un patrón.
\end{itemize}

\subsection{Datos: cinco años de ventas mensuales}

Las ventas del curso solo cubren 2023 y son aleatorias, así que simulamos 5
años (2021--2025) de unidades vendidas por mes con una tendencia creciente
(unas 15 unidades más cada mes), una estacionalidad típica del retail mexicano
(picos de 15\% y 32\% sobre lo normal en noviembre y diciembre) y 3\% de ruido.

\begin{rcode}
set.seed(123)
meses <- seq(as.Date("2021-01-01"), as.Date("2025-12-01"), by = "month")
n_meses <- length(meses)   # 60 meses

# Factor estacional de cada mes (1 = mes normal)
factor_mes <- c(0.88, 0.85, 0.95, 0.97, 1.04, 0.96,   # ene-jun
                0.97, 1.00, 0.93, 0.98, 1.15, 1.32)   # jul-dic
factor_mes <- factor_mes / mean(factor_mes)            # promedio = 1

ventas_mensuales <- tibble(
  mes = meses,
  tendencia = 1500 + 15 * (1:n_meses),
  unidades = round(tendencia * factor_mes[month(meses)] *
                     (1 + rnorm(n_meses, 0, 0.03)))    # ruido de 3 %
) %>%
  select(-tendencia)

# Objeto ts: la estructura de series de tiempo de R
serie <- ts(ventas_mensuales$unidades, start = c(2021, 1),
            frequency = 12)   # 12 = datos mensuales
serie
\end{rcode}

\begin{rsalida}
      Jan  Feb  Mar  Apr  May  Jun  Jul  Aug  Sep  Oct  Nov  Dec
2021 1311 1292 1536 1516 1644 1605 1578 1559 1489 1595 1985 2242
2022 1510 1458 1611 1778 1852 1599 1768 1774 1634 1782 2056 2402
2023 1619 1525 1855 1871 1944 1942 1930 1962 1905 2022 2386 2748
2024 1838 1756 1963 2014 2154 2032 2002 2301 2096 2074 2505 2889
2025 2013 1908 2168 2210 2384 2309 2240 2446 2088 2363 2753 3189
\end{rsalida}

El objeto \codigo{ts} guarda los valores junto con su calendario: inicio
(\codigo{start = c(2021, 1)}, enero de 2021) y frecuencia (12 observaciones por
año; usa 4 para trimestres, 52 para semanas). Por eso R imprime la serie como
una tabla año por mes (los nombres de los meses siempre salen en inglés en este
formato). Ya a simple vista se nota el crecimiento de un año a
otro y los picos de diciembre.

\begin{rcode}
grafico_serie <- autoplot(serie, color = color_principal,
                          linewidth = 0.8) +
  scale_y_continuous(labels = scales::comma) +
  labs(title = "Unidades vendidas por mes, electrónica",
       subtitle = "2021-2025 (datos simulados)",
       x = NULL, y = "Unidades") +
  theme_minimal() +
  theme(panel.grid.minor = element_blank())
grafico_serie
\end{rcode}

\begin{figure}[H]
\centering
\includegraphics[width=0.9\textwidth]{m07-serie}
\caption{Serie mensual de unidades vendidas: tendencia creciente y picos cada
fin de año.}
\label{fig:m07-serie}
\end{figure}

\codigo{autoplot()} de \paquete{forecast} convierte objetos de series de tiempo
en gráficas de \paquete{ggplot2}, así que puedes agregarles capas, escalas y
temas como a cualquier otra.

\subsection{Descomposición: separar tendencia, estacionalidad y ruido}

\begin{rcode}
descomposicion <- decompose(serie, type = "multiplicative")

# Índices estacionales: cuánto se aleja cada mes de un mes "normal"
nombres_mes <- c("ene", "feb", "mar", "abr", "may", "jun",
                 "jul", "ago", "sep", "oct", "nov", "dic")
tibble(mes = nombres_mes, indice = round(descomposicion$figure, 3)) %>%
  arrange(desc(indice)) %>%
  head(4)

grafico_descomp <- autoplot(descomposicion) +
  labs(title = "Descomposición multiplicativa de la serie",
       x = NULL) +
  theme_minimal() +
  theme(panel.grid.minor = element_blank())
grafico_descomp
\end{rcode}

\begin{rsalida}
# A tibble: 4 x 2
  mes   indice
  <chr>  <dbl>
1 dic     1.32
2 nov     1.16
3 may     1.03
4 ago     1   
\end{rsalida}

\begin{figure}[H]
\centering
\includegraphics[width=0.9\textwidth]{m07-descomposicion}
\caption{Descomposición de la serie: datos originales, tendencia, componente
estacional y residuo (\emph{remainder}).}
\label{fig:m07-descomposicion}
\end{figure}

Usamos el modelo \emph{multiplicativo} porque la estacionalidad es un
porcentaje (diciembre vende ``32\% más'', no ``400 unidades más''), y ese
porcentaje se mantiene aunque el negocio crezca. Los índices estacionales recuperan
casi exactamente lo que simulamos: diciembre vende 32\% más que un mes normal
(índice 1.32) y noviembre 16\% más (1.16). Esos índices, por sí solos, ya son
información valiosa para planear personal, inventario y presupuesto de
publicidad.

\begin{hazlotu}
Usa \codigo{stl(serie, s.window = "periodic")} en lugar de \codigo{decompose()}
y grafica el resultado con \codigo{autoplot()}. STL es un método de
descomposición más robusto ante valores atípicos (trabaja de forma aditiva;
para una versión multiplicativa aplica \codigo{log()} a la serie).
\end{hazlotu}

\subsection{Entrenamiento y prueba en series de tiempo}

En series de tiempo \textbf{no} se separan los datos al azar: el pasado entrena
y el futuro evalúa. Guardamos los últimos 12 meses (todo 2025) como prueba, que
es exactamente lo que queremos saber hacer: pronosticar un año.

\begin{rcode}
serie_ent    <- window(serie, end = c(2024, 12))   # 2021-2024: 48 meses
serie_prueba <- window(serie, start = c(2025, 1))  # 2025: 12 meses
c(length(serie_ent), length(serie_prueba))
\end{rcode}

\begin{rsalida}
[1] 48 12
\end{rsalida}

\subsection{Cuatro modelos de pronóstico}

\begin{concepto}{Modelos de pronóstico}
\begin{itemize}
  \item \textbf{Ingenuo} (\codigo{naive()}): el próximo mes será igual al
        último. Parece tonto, pero es la línea base obligatoria.
  \item \textbf{Ingenuo estacional} (\codigo{snaive()}): cada mes será igual
        al mismo mes del año anterior (diciembre 2025 = diciembre 2024).
  \item \textbf{ETS} (\codigo{ets()}, \emph{Error, Trend, Seasonal}):
        suavizamiento exponencial; da más peso a los datos recientes y modela
        explícitamente nivel, tendencia y estacionalidad. Elige solo la mejor
        combinación.
  \item \textbf{ARIMA} (\codigo{auto.arima()}): modela la relación de cada mes
        con los meses anteriores y con sus propios errores pasados;
        \codigo{auto.arima()} busca la mejor especificación automáticamente.
\end{itemize}
\end{concepto}

\begin{rcode}
h <- length(serie_prueba)   # horizonte: 12 meses

pron_naive  <- naive(serie_ent, h = h)
pron_snaive <- snaive(serie_ent, h = h)
modelo_ets  <- ets(serie_ent)
pron_ets    <- forecast(modelo_ets, h = h)
modelo_arima <- auto.arima(serie_ent)
pron_arima  <- forecast(modelo_arima, h = h)

modelo_ets     # ¿qué combinación eligió ets()?
modelo_arima
\end{rcode}

\begin{rsalida}
ETS(M,A,M) 

Call:
 ets(y = serie_ent) 

  Smoothing parameters:
    alpha = 0.0199 
    beta  = 4e-04 
    gamma = 1e-04 

  Initial states:
    l = 1527.3166 
    b = 14.5543 
    s = 1.3231 1.1584 0.9739 0.9307 1.0039 0.9739
           0.9589 1.0257 0.9811 0.9589 0.8368 0.8747

  sigma:  0.0339

     AIC     AICc      BIC 
597.8241 618.2241 629.6346 
Series: serie_ent 
ARIMA(0,0,0)(1,1,0)[12] with drift 

Coefficients:
         sar1    drift
      -0.5378  14.9348
s.e.   0.1418   0.7466

sigma^2 = 5463:  log likelihood = -207
AIC=420.01   AICc=420.76   BIC=424.76
\end{rsalida}

\codigo{ets()} eligió un modelo \codigo{ETS(M,A,M)}: error multiplicativo,
tendencia aditiva (\emph{A}, crece una cantidad fija cada mes: \codigo{b =
14.55}, muy cerca de las 15 unidades que simulamos) y estacionalidad
multiplicativa (\emph{M}). \codigo{auto.arima()} eligió un modelo que compara
cada mes con el mismo mes del año anterior (la parte \codigo{(1,1,0)[12]}) más
una deriva (\emph{drift}) de 14.9 unidades por mes: otra forma de describir la
misma tendencia. No necesitas memorizar la notación; lo importante es que
ambos modelos capturaron tendencia y estacionalidad sin que se las
indicáramos.

\subsection{Comparar los modelos con \texorpdfstring{\codigo{accuracy()}}{accuracy()}}

\codigo{accuracy()} calcula las métricas de error del pronóstico contra los
valores reales. Nos interesa la fila \codigo{Test set} (error en 2025, que los
modelos no vieron):

\begin{rcode}
accuracy(pron_ets, serie_prueba)   # salida completa de un modelo

comparacion_ts <- list(Ingenuo = pron_naive,
                       `Ingenuo estacional` = pron_snaive,
                       ETS = pron_ets, ARIMA = pron_arima) %>%
  map_dfr(function(p) {
    acc <- accuracy(p, serie_prueba)["Test set", ]
    tibble(MAE = acc["MAE"], RMSE = acc["RMSE"], MAPE = acc["MAPE"])
  }, .id = "modelo") %>%
  mutate(across(where(is.numeric), ~ round(.x, 1))) %>%
  arrange(MAPE)
comparacion_ts
\end{rcode}

\begin{rsalida}
                    ME     RMSE      MAE        MPE     MAPE
Training set -7.685706 52.06332 43.34896 -0.5184226 2.336622
Test set     27.712908 60.48803 48.44893  1.1527656 2.130479
                  MASE       ACF1 Theil's U
Training set 0.2483390  0.1141467        NA
Test set     0.2775559 -0.5886038 0.2412182
# A tibble: 4 x 4
  modelo               MAE  RMSE  MAPE
  <chr>              <dbl> <dbl> <dbl>
1 ARIMA               46.8  68.3   2  
2 ETS                 48.4  60.5   2.1
3 Ingenuo estacional 205.  219.    8.7
4 Ingenuo            600.  642    27.3
\end{rsalida}

La salida completa de \codigo{accuracy()} incluye más métricas (\codigo{ME}
es el error promedio con signo: 27.7 indica que ETS se quedó ligeramente corto
en 2025; \codigo{MASE} compara contra el ingenuo). En la tabla resumida, ETS y
ARIMA prácticamente empatan con errores de apenas 2\% (MAPE) en 2025: ARIMA
tiene un MAE algo menor (46.8 contra 48.4) y ETS un RMSE menor (60.5 contra
68.3), es decir, menos errores grandes. El ingenuo estacional se equivoca 8.7\%
porque ignora el crecimiento, y el ingenuo simple, 27.3\%, porque ignora la
estacionalidad: diciembre de 2024 no se parece a enero de 2025. Elegimos ETS
por su menor RMSE: en inventario, los errores grandes son los que duelen.

\begin{rcode}
grafico_comparacion <- autoplot(window(serie, start = c(2023, 1)),
                                color = "black", linewidth = 0.6) +
  autolayer(pron_snaive, series = "Ingenuo estacional", PI = FALSE) +
  autolayer(pron_ets, series = "ETS", PI = FALSE) +
  autolayer(pron_arima, series = "ARIMA", PI = FALSE) +
  scale_color_manual(values = paleta_libro[c(2, 1, 3)]) +
  scale_y_continuous(labels = scales::comma) +
  labs(title = "Pronósticos de 2025 contra la realidad",
       subtitle = "Negro: ventas reales; entrenamiento hasta dic-2024",
       x = NULL, y = "Unidades", color = NULL) +
  theme_minimal() +
  theme(panel.grid.minor = element_blank(), legend.position = "top")
grafico_comparacion
\end{rcode}

\begin{figure}[H]
\centering
\includegraphics[width=0.9\textwidth]{m07-comparacion-pronosticos}
\caption{Pronósticos de 2025 de tres modelos contra las ventas reales. ETS y
ARIMA siguen de cerca la realidad; el ingenuo estacional se queda corto porque
ignora el crecimiento.}
\label{fig:m07-comparacion-pronosticos}
\end{figure}

\subsection{Pronóstico final con intervalos}

Una vez elegido el modelo, se \textbf{reentrena con todos los datos}
(incluido 2025) para pronosticar 2026:

\begin{rcode}
modelo_final <- ets(serie)
pronostico_2026 <- forecast(modelo_final, h = 12, level = c(80, 95))
pronostico_2026
\end{rcode}

\begin{rsalida}
         Point Forecast    Lo 80    Hi 80    Lo 95    Hi 95
Jan 2026       2125.663 2040.343 2210.984 1995.177 2256.150
Feb 2026       2037.570 1955.751 2119.389 1912.439 2162.702
Mar 2026       2333.729 2239.979 2427.480 2190.350 2477.108
Apr 2026       2400.133 2303.674 2496.591 2252.612 2547.653
May 2026       2547.002 2444.599 2649.406 2390.390 2703.615
Jun 2026       2411.232 2314.247 2508.217 2262.906 2559.557
Jul 2026       2420.791 2323.381 2518.200 2271.816 2569.766
Aug 2026       2548.945 2446.337 2651.553 2392.020 2705.871
Sep 2026       2349.700 2255.074 2444.327 2204.982 2494.419
Oct 2026       2488.212 2387.967 2588.457 2334.901 2641.523
Nov 2026       2975.194 2855.281 3095.107 2791.803 3158.585
Dec 2026       3411.931 3274.360 3549.501 3201.535 3622.327
\end{rsalida}

\begin{concepto}{Intervalos de predicción}
\codigo{Point Forecast} es la estimación central. \codigo{Lo 80} y \codigo{Hi 80}
forman el intervalo del 80\%: esperamos que 8 de cada 10 meses la venta real
caiga dentro de ese rango. El intervalo del 95\% es más ancho porque exige más
seguridad. Los intervalos se ensanchan a medida que el pronóstico se aleja:
es más difícil predecir diciembre que enero.
\end{concepto}

\begin{rcode}
grafico_pronostico <- autoplot(pronostico_2026, color = color_principal,
                               fcol = color_resalte, flwd = 0.8) +
  scale_y_continuous(labels = scales::comma) +
  labs(title = "Pronóstico de unidades para 2026",
       subtitle = "Modelo ETS; bandas de 80 % y 95 %",
       x = NULL, y = "Unidades") +
  theme_minimal() +
  theme(panel.grid.minor = element_blank())
grafico_pronostico
\end{rcode}

\begin{figure}[H]
\centering
\includegraphics[width=0.9\textwidth]{m07-pronostico}
\caption{Pronóstico de 12 meses con intervalos de predicción del 80\% (banda
oscura) y 95\% (banda clara).}
\label{fig:m07-pronostico}
\end{figure}

\subsection{Comunicar la incertidumbre a la dirección}

Un pronóstico sin rango es una promesa que no se va a cumplir exactamente.
Presenta siempre tres escenarios en lenguaje de negocio:

\begin{rcode}
tabla_direccion <- tibble(
  mes         = format(seq(as.Date("2026-01-01"), by = "month",
                           length.out = 12), "%b %Y"),
  pesimista   = round(pronostico_2026$lower[, "80%"]),
  esperado    = round(pronostico_2026$mean),
  optimista   = round(pronostico_2026$upper[, "80%"])
)
tabla_direccion %>% slice(c(1, 6, 11, 12))
sum(tabla_direccion$esperado)   # unidades esperadas en todo 2026
\end{rcode}

\begin{rsalida}
# A tibble: 4 x 4
  mes      pesimista esperado optimista
  <chr>        <dbl>    <dbl>     <dbl>
1 ene 2026      2040     2126      2211
2 jun 2026      2314     2411      2508
3 nov 2026      2855     2975      3095
4 dic 2026      3274     3412      3550
[1] 30051
\end{rsalida}

Para la dirección: en diciembre de 2026 esperamos vender 3{,}412 unidades; en
un escenario pesimista (límite inferior del 80\%), 3{,}274, y en uno optimista,
3{,}550. En todo el año se esperan 30{,}051 unidades. Observa que
\codigo{format(..., "\%b \%Y")} sí respeta el idioma del sistema (``ene
2026'').

\begin{negocio}{Cómo presentar un pronóstico}
\begin{itemize}
  \item Da un rango y un escenario central: ``esperamos vender 3{,}412 unidades
        en diciembre; con 80\% de confianza, entre 3{,}274 y 3{,}550''.
  \item Explica los supuestos: ``suponemos que el crecimiento y la
        estacionalidad de los últimos cinco años se mantienen''.
  \item Di qué lo podría romper: una devaluación, un competidor nuevo, una
        promoción extraordinaria, una pandemia.
  \item Compara con el pronóstico anterior y con lo que realmente pasó: la
        credibilidad se gana mostrando el historial de aciertos y errores.
\end{itemize}
\end{negocio}

\subsection{Del pronóstico al inventario: stock de seguridad}

El pronóstico central no basta para decidir cuánto pedir: si pides exactamente
lo esperado, la mitad de los meses te faltará mercancía. El
\termino{stock de seguridad} es el inventario extra que cubre la
incertidumbre. Una fórmula sencilla y muy usada es:

\[
\text{stock de seguridad} = z \times \sigma_{\text{error}} \times \sqrt{L}
\]

donde $z$ depende del \termino{nivel de servicio} deseado (1.28 para 90\%,
1.65 para 95\%, 2.33 para 99\%), $\sigma_{\text{error}}$ es la desviación
estándar del error de pronóstico (usaremos el RMSE del modelo ETS en 2025) y
$L$ es el tiempo de reabastecimiento en meses.

\begin{rcode}
nivel_servicio <- 0.95
z <- qnorm(nivel_servicio)                 # 1.645
error_tipico <- comparacion_ts %>%
  filter(modelo == "ETS") %>% pull(RMSE)   # error típico en 2025
tiempo_entrega <- 2                        # meses de anticipación

stock_seguridad <- z * error_tipico * sqrt(tiempo_entrega)

tabla_direccion %>%
  slice(11:12) %>%
  mutate(stock_seguridad = round(stock_seguridad),
         pedido_sugerido = esperado + stock_seguridad) %>%
  select(mes, esperado, stock_seguridad, pedido_sugerido)
\end{rcode}

\begin{rsalida}
# A tibble: 2 x 4
  mes      esperado stock_seguridad pedido_sugerido
  <chr>       <dbl>           <dbl>           <dbl>
1 nov 2026     2975             141            3116
2 dic 2026     3412             141            3553
\end{rsalida}

Con un nivel de servicio de 95\% ($z = 1.645$), un error típico de 60.5
unidades y dos meses de anticipación, el stock de seguridad es de 141
unidades. El pedido sugerido para diciembre es de 3{,}553 unidades: el
pronóstico central más el colchón. Subir el nivel de servicio a 99\% aumenta el
colchón (y el costo de inventario); es otra decisión de negocio que el
modelo te ayuda a cuantificar.

\begin{advertencia}
Un modelo de series de tiempo asume que \textbf{el futuro se parece al
pasado}. No sabe que el próximo año abrirás 20 tiendas, que un proveedor
quebró o que habrá una promoción extraordinaria. Pronostica a 12 meses; más
allá, la incertidumbre crece tanto que el número deja de ser útil para
decisiones operativas. Y revisa el pronóstico cada mes con los datos nuevos.
\end{advertencia}

% ============================================================================
\section{Caso 5: análisis de canasta (market basket)}
\label{sec:m7-canasta}

\begin{ejemplo}{¿Qué le ofrecemos al cliente que compra una laptop?}
Una tienda de electrónica y papelería quiere diseñar promociones de venta
cruzada (\emph{cross-selling}): ``lleva el mouse con 20\% de descuento'',
colocar productos juntos en el anaquel o recomendar productos en la tienda en
línea. La pregunta: \emph{¿qué productos se compran juntos con más frecuencia
de la que esperaríamos por casualidad?}
\end{ejemplo}

El \termino{análisis de canasta} (\emph{market basket analysis}) busca
\termino{reglas de asociación} del tipo \{Laptop\} $\rightarrow$ \{Mouse\}:
``quien compra una laptop tiende a comprar un mouse''. Es aprendizaje no
supervisado: no hay nada que predecir, solo patrones que descubrir en los
tickets de compra.

\subsection{Soporte, confianza y lift}

Tres medidas describen una regla $X \rightarrow Y$:

\begin{itemize}
  \item \termino{Soporte}: proporción de tickets que contienen $X$ \emph{y}
        $Y$. Mide qué tan frecuente es la combinación.
        \[ \text{soporte}(X \rightarrow Y) = \frac{\text{tickets con } X \text{ y } Y}{\text{total de tickets}} \]
  \item \termino{Confianza}: de los tickets que tienen $X$, qué proporción
        también tiene $Y$. Es una probabilidad condicional.
        \[ \text{confianza}(X \rightarrow Y) = \frac{\text{tickets con } X \text{ y } Y}{\text{tickets con } X} \]
  \item \termino{Lift}: cuántas veces más probable es comprar $Y$ cuando se
        compra $X$, comparado con comprar $Y$ en general.
        \[ \text{lift}(X \rightarrow Y) = \frac{\text{confianza}(X \rightarrow Y)}{\text{soporte}(Y)} \]
        Lift $= 1$: no hay relación; lift $> 1$: se compran juntos más de lo
        esperado; lift $< 1$: se ``evitan''.
\end{itemize}

\textbf{Ejemplo a mano.} En 1{,}000 tickets, 100 tienen laptop, 200 tienen
mouse y 60 tienen ambos.

\begin{itemize}
  \item Soporte de \{Laptop, Mouse\}: $60 / 1{,}000 = 0.06$ (6\% de los tickets).
  \item Confianza de Laptop $\rightarrow$ Mouse: $60 / 100 = 0.60$ (60\% de
        quienes compran laptop llevan mouse).
  \item Soporte de Mouse: $200 / 1{,}000 = 0.20$ (20\% de todos los tickets
        llevan mouse).
  \item Lift: $0.60 / 0.20 = 3$. Quien compra laptop tiene \textbf{tres veces}
        más probabilidad de llevar mouse que un cliente cualquiera.
\end{itemize}

\begin{advertencia}
La confianza sola engaña. Si el 30\% de todos los tickets lleva papel, una
regla ``Laptop $\rightarrow$ Papel'' con confianza de 30\% parece buena, pero
su lift es 1: el papel se vende igual con o sin laptop. Por eso ordenamos las
reglas por lift y exigimos un soporte mínimo (las reglas con muy pocos tickets
son casualidades).
\end{advertencia}

\subsection{¿Sirven los datos del curso?}

Antes de simular, revisemos si \archivo{transacciones.csv} permite formar
canastas. Consideremos una canasta como todos los productos que un cliente
compró el mismo día:

\begin{rcode}
canastas_curso <- transacciones %>%
  group_by(cliente_id, fecha) %>%
  summarise(productos = n_distinct(producto_id), .groups = "drop")

canastas_curso %>% count(productos, name = "canastas")
\end{rcode}

\begin{rsalida}
# A tibble: 2 x 2
  productos canastas
      <int>    <int>
1         1      974
2         2       13
\end{rsalida}

De 987 canastas, solo 13 tienen dos productos; ninguna tiene más. Con los datos
del curso no hay canastas que analizar (cada transacción es prácticamente una
compra de un solo producto). Otra vez: sin señal no hay modelo. Simulamos, por
lo tanto, 3{,}000 tickets de una tienda con asociaciones conocidas.

\subsection{Tickets simulados con asociaciones conocidas}

Ocho productos se compran de forma independiente, y cuatro dependen de otro:
el mouse y la funda son mucho más probables si el ticket incluye una laptop, el
cable HDMI si incluye un monitor y la tinta si incluye una impresora.

\begin{rcode}
set.seed(123)
n_tickets <- 3000
tickets_ancho <- tibble(
  ticket_id = 1:n_tickets,
  Laptop    = rbinom(n_tickets, 1, 0.08),   # 1 = el ticket lo incluye
  Monitor   = rbinom(n_tickets, 1, 0.07),
  Impresora = rbinom(n_tickets, 1, 0.05),
  Audifonos = rbinom(n_tickets, 1, 0.15),
  USB       = rbinom(n_tickets, 1, 0.18),
  Papel     = rbinom(n_tickets, 1, 0.20),
  Teclado   = rbinom(n_tickets, 1, 0.08)
) %>%
  mutate(
    # Asociaciones "reales" que el análisis debe descubrir
    Mouse = if_else(Laptop == 1, rbinom(n(), 1, 0.60),
                    rbinom(n(), 1, 0.08)),
    Funda = if_else(Laptop == 1, rbinom(n(), 1, 0.40),
                    rbinom(n(), 1, 0.02)),
    `Cable HDMI` = if_else(Monitor == 1, rbinom(n(), 1, 0.50),
                           rbinom(n(), 1, 0.04)),
    Tinta = if_else(Impresora == 1, rbinom(n(), 1, 0.70),
                    rbinom(n(), 1, 0.06))
  )

# Formato largo: una fila por ticket y producto (como en un sistema real)
tickets <- tickets_ancho %>%
  pivot_longer(-ticket_id, names_to = "producto", values_to = "compro") %>%
  filter(compro == 1) %>%
  select(-compro)

head(tickets, 5)
n_total <- n_distinct(tickets$ticket_id)   # tickets con al menos 1 producto
n_total
\end{rcode}

\begin{rsalida}
# A tibble: 5 x 2
  ticket_id producto 
      <int> <chr>    
1         1 Teclado  
2         2 Monitor  
3         2 Audifonos
4         3 USB      
5         4 USB      
[1] 1975
\end{rsalida}

El formato largo (una fila por ticket y producto) es exactamente como vendrían
los datos de un sistema de punto de venta. Hay 1{,}975 tickets con al menos un
producto (los tickets vacíos de la simulación se descartan).

\subsection{Calcular reglas de pares con \texorpdfstring{\paquete{dplyr}}{dplyr}}

Primero el soporte de cada producto:

\begin{rcode}
soporte_producto <- tickets %>%
  count(producto, name = "tickets_x") %>%
  mutate(soporte_x = tickets_x / n_total) %>%
  arrange(desc(soporte_x))
soporte_producto %>% mutate(soporte_x = round(soporte_x, 3))
\end{rcode}

\begin{rsalida}
# A tibble: 11 x 3
   producto   tickets_x soporte_x
   <chr>          <int>     <dbl>
 1 Papel            596     0.302
 2 USB              557     0.282
 3 Audifonos        442     0.224
 4 Mouse            386     0.195
 5 Tinta            276     0.14 
 6 Laptop           225     0.114
 7 Cable HDMI       215     0.109
 8 Teclado          210     0.106
 9 Monitor          205     0.104
10 Impresora        149     0.075
11 Funda            144     0.073
\end{rsalida}

Ahora el truco: un \termino{self-join} (unir la tabla consigo misma por
\codigo{ticket\_id}) genera todos los pares de productos que aparecen en el
mismo ticket. Luego contamos cuántos tickets tiene cada par y calculamos las
tres medidas:

\begin{rcode}
reglas <- tickets %>%
  inner_join(tickets, by = "ticket_id", suffix = c("_x", "_y"),
             relationship = "many-to-many") %>%
  filter(producto_x != producto_y) %>%          # sin pares consigo mismo
  count(producto_x, producto_y, name = "tickets_xy") %>%
  left_join(soporte_producto, by = c("producto_x" = "producto")) %>%
  left_join(soporte_producto %>%
              select(producto_y = producto, soporte_y = soporte_x),
            by = "producto_y") %>%
  mutate(soporte   = tickets_xy / n_total,
         confianza = tickets_xy / tickets_x,
         lift      = confianza / soporte_y)

mejores_reglas <- reglas %>%
  filter(soporte >= 0.02) %>%                   # soporte mínimo: 2 %
  arrange(desc(lift)) %>%
  transmute(regla = paste(producto_x, "->", producto_y),
            tickets_xy, soporte = round(soporte, 3),
            confianza = round(confianza, 3), lift = round(lift, 2))
mejores_reglas %>% head(10)
\end{rcode}

\begin{rsalida}
# A tibble: 10 x 5
   regla                 tickets_xy soporte confianza  lift
   <chr>                      <int>   <dbl>     <dbl> <dbl>
 1 Funda -> Laptop               82   0.042     0.569  5   
 2 Laptop -> Funda               82   0.042     0.364  5   
 3 Impresora -> Tinta           104   0.053     0.698  4.99
 4 Tinta -> Impresora           104   0.053     0.377  4.99
 5 Cable HDMI -> Monitor        105   0.053     0.488  4.71
 6 Monitor -> Cable HDMI        105   0.053     0.512  4.71
 7 Laptop -> Mouse              136   0.069     0.604  3.09
 8 Mouse -> Laptop              136   0.069     0.352  3.09
 9 Funda -> Mouse                61   0.031     0.424  2.17
10 Mouse -> Funda                61   0.031     0.158  2.17
\end{rsalida}

El análisis descubrió exactamente las cuatro asociaciones que simulamos, y en
ambas direcciones:

\begin{itemize}
  \item \textbf{Impresora $\rightarrow$ Tinta}: confianza de 69.8\% (simulamos
        70\%) y lift de 4.99. Casi 7 de cada 10 compradores de impresora llevan
        tinta.
  \item \textbf{Monitor $\rightarrow$ Cable HDMI}: confianza de 51.2\% (simulamos
        50\%), lift de 4.71.
  \item \textbf{Laptop $\rightarrow$ Mouse}: confianza de 60.4\% (simulamos
        60\%), lift de 3.09, y es la regla con más tickets (136).
  \item \textbf{Laptop $\rightarrow$ Funda}: lift de 5.0, el más alto.
\end{itemize}

Observa que el lift es simétrico (Laptop $\rightarrow$ Funda y Funda
$\rightarrow$ Laptop tienen el mismo), pero la confianza no: 56.9\% de quienes
compran funda llevan laptop, pero solo 36.4\% de quienes compran laptop llevan
funda. Y aparece una regla que \emph{no} simulamos directamente: Funda
$\rightarrow$ Mouse (lift 2.17). No es que la funda ``cause'' el mouse: ambos
se compran con laptop. Las reglas de asociación describen coocurrencias, no
causas.

\begin{rcode}
grafico_lift <- mejores_reglas %>%
  filter(lift > 1.5) %>%
  ggplot(aes(x = lift, y = reorder(regla, lift))) +
  geom_col(fill = color_principal, width = 0.7) +
  geom_vline(xintercept = 1, color = color_gris, linetype = "dashed") +
  labs(title = "Reglas de asociación con mayor lift",
       subtitle = "Soporte mínimo 2 %; la línea punteada es lift = 1",
       x = "Lift (veces más probable que al azar)", y = NULL) +
  theme_minimal() +
  theme(panel.grid.minor = element_blank(),
        panel.grid.major.y = element_blank())
grafico_lift
\end{rcode}

\begin{figure}[H]
\centering
\includegraphics[width=0.9\textwidth]{m07-lift}
\caption{Reglas de asociación ordenadas por lift. Las diez reglas corresponden
a los cuatro pares que simulamos, en ambas direcciones, más la relación
indirecta Funda--Mouse.}
\label{fig:m07-lift}
\end{figure}

¿Y las reglas con lift menor a 1? También dan información:

\begin{rcode}
mejores_reglas %>% arrange(lift) %>% head(3)
\end{rcode}

\begin{rsalida}
# A tibble: 3 x 5
  regla              tickets_xy soporte confianza  lift
  <chr>                   <int>   <dbl>     <dbl> <dbl>
1 Papel -> USB              100   0.051     0.168  0.59
2 USB -> Papel              100   0.051     0.18   0.59
3 Audifonos -> Mouse         52   0.026     0.118  0.6 
\end{rsalida}

Papel y USB aparecen juntos \emph{menos} de lo esperado (lift 0.59), aunque
los simulamos de forma independiente. Es un efecto de cómo se forman los datos:
al quedarnos solo con tickets que tienen al menos un producto, los productos
frecuentes ``compiten'' por aparecer. En datos reales, un lift claramente
menor a 1 suele indicar \textbf{productos sustitutos} (dos marcas de la misma
tinta, por ejemplo): no tiene sentido promocionarlos juntos. La lección es que
ningún número se interpreta sin entender cómo se generaron los datos.

\subsection{Con el paquete arules}

Para canastas con cientos de productos y reglas de más de dos productos
(\{Laptop, Mouse\} $\rightarrow$ \{Funda\}), el estándar en R es el paquete
\paquete{arules}, que implementa el algoritmo \termino{Apriori}. No está
instalado en este entorno, pero así se usaría con los mismos datos:

\begin{rnoejecutar}
install.packages("arules")
library(arules)

# Convertir a formato de transacciones: una lista de productos por ticket
lista_tickets <- split(tickets$producto, tickets$ticket_id)
trans <- as(lista_tickets, "transactions")
summary(trans)

reglas_apriori <- apriori(trans,
                          parameter = list(supp = 0.02, conf = 0.3,
                                           minlen = 2))
inspect(head(sort(reglas_apriori, by = "lift"), 10))
\end{rnoejecutar}

\subsection{De las reglas a las recomendaciones}

\begin{negocio}{Acciones de venta cruzada}
\begin{itemize}
  \item \textbf{Paquetes}: ``Kit oficina en casa'' con Impresora + Tinta, o
        ``Kit estudiante'' con Laptop + Mouse + Funda, con un pequeño descuento
        sobre la suma.
  \item \textbf{Ubicación en tienda}: los cables HDMI junto a los monitores;
        las tintas a la vista desde el pasillo de impresoras.
  \item \textbf{Recomendaciones en línea}: al agregar una laptop al carrito,
        mostrar mouse y funda (las reglas con mayor confianza desde Laptop).
  \item \textbf{No descontar lo que ya se vende junto}: si el 70\% de quienes
        compran impresora ya lleva tinta, un descuento en tinta regala
        margen. Mejor ofrecer el descuento al 30\% restante (por ejemplo, en
        el correo de seguimiento).
  \item \textbf{Medir}: lanza la promoción en algunas tiendas y compara con
        otras similares sin promoción antes de extenderla.
\end{itemize}
\end{negocio}

\begin{hazlotu}
Calcula las reglas cuyo antecedente es \codigo{"Laptop"}
(\codigo{filter(producto\_x == "Laptop")}) y ordénalas por confianza. ¿Qué
producto recomendarías primero en el carrito de compras? ¿Coincide con el orden
por lift?
\end{hazlotu}

% ============================================================================
\section{Llevar el modelo a producción}
\label{sec:m7-produccion}

Un modelo que vive solo en tu computadora no genera valor. \termino{Desplegar}
(o ``poner en producción'') un modelo significa que se ejecuta de forma
periódica y confiable sobre datos nuevos, y que sus resultados llegan a quien
toma decisiones. En BI, lo más común es empezar con algo sencillo: un script
que cada lunes puntúa a los clientes y deja un CSV (o una tabla en la base de
datos, Módulo~\ref{cap:bases-datos}) que alimenta el tablero y el call center.

\subsection{Guardar y cargar el modelo}

No necesitas reentrenar cada vez: guarda el objeto del modelo en un archivo
\archivo{.rds} con \codigo{saveRDS()} y cárgalo con \codigo{readRDS()}.

\begin{rcode}
dir.create("resultados", showWarnings = FALSE)

saveRDS(modelo_log, "resultados/modelo_churn.rds")

# ... días después, en otro script o en otra computadora:
modelo_cargado <- readRDS("resultados/modelo_churn.rds")

# Comprobamos que predice exactamente lo mismo
all.equal(predict(modelo_cargado, churn_prueba, type = "response"),
          predict(modelo_log, churn_prueba, type = "response"))
\end{rcode}

\begin{rsalida}
[1] TRUE
\end{rsalida}

\begin{consejo}
Guarda junto al modelo la información para reproducirlo: fecha de
entrenamiento, periodo de datos, métricas en prueba y versión de R y de los
paquetes. Un nombre como \archivo{modelo\_churn\_2025-06\_v2.rds} y una nota
en un archivo de texto te ahorrarán muchos dolores de cabeza.
\end{consejo}

\subsection{Una función para puntuar clientes nuevos}

Encapsula la predicción en una función que \textbf{valide} los datos de
entrada, calcule la probabilidad y la traduzca a una acción:

\begin{rcode}
predecir_churn <- function(nuevos_datos, modelo, umbral = 0.15) {
  # 1. Validar que vengan todas las columnas que el modelo necesita
  necesarias <- c("antiguedad_meses", "plan", "uso_mensual", "quejas_6m",
                  "pagos_tardios", "genero", "edad")
  faltan <- setdiff(necesarias, names(nuevos_datos))
  if (length(faltan) > 0) {
    stop("Faltan columnas: ", paste(faltan, collapse = ", "))
  }
  # 2. Mismos niveles de factor que en el entrenamiento
  nuevos_datos <- nuevos_datos %>%
    mutate(plan = factor(plan, levels = c("Básico", "Estándar",
                                          "Premium")))
  # 3. Probabilidad y acción sugerida
  nuevos_datos %>%
    mutate(
      prob_churn = predict(modelo, newdata = nuevos_datos,
                           type = "response"),
      nivel_riesgo = cut(prob_churn, breaks = c(0, umbral, 0.5, 1),
                         labels = c("Bajo", "Medio", "Alto"),
                         include.lowest = TRUE),
      accion = if_else(prob_churn >= umbral, "Llamar", "Sin acción")
    )
}
\end{rcode}

Simulamos un lote de 500 clientes activos de esta semana (en la vida real
vendría de la base de datos) y lo puntuamos:

\begin{rcode}
set.seed(123)
lote_semana <- tibble(
  cliente_id       = sprintf("N%04d", 1:500),
  antiguedad_meses = sample(1:72, 500, replace = TRUE),
  plan             = sample(c("Básico", "Estándar", "Premium"), 500,
                            replace = TRUE, prob = c(0.5, 0.3, 0.2)),
  uso_mensual      = rpois(500, 8),
  quejas_6m        = rpois(500, 0.8),
  pagos_tardios    = rpois(500, 0.4),
  genero           = sample(c("F", "M"), 500, replace = TRUE),
  edad             = sample(18:70, 500, replace = TRUE)
)

puntuados <- predecir_churn(lote_semana, modelo_cargado)
puntuados %>% count(nivel_riesgo, accion)

# Exportar la lista para el equipo de retención
puntuados %>%
  filter(accion == "Llamar") %>%
  arrange(desc(prob_churn)) %>%
  mutate(prob_churn = round(prob_churn, 3)) %>%
  select(cliente_id, plan, quejas_6m, prob_churn, nivel_riesgo) %>%
  write_csv("resultados/clientes_a_llamar.csv")

file.exists("resultados/clientes_a_llamar.csv")
\end{rcode}

\begin{rsalida}
# A tibble: 3 x 3
  nivel_riesgo accion         n
  <fct>        <chr>      <int>
1 Bajo         Sin acción   275
2 Medio        Llamar       142
3 Alto         Llamar        83
[1] TRUE
\end{rsalida}

De los 500 clientes, 225 superan el umbral de 0.15 y se llaman: 83 de riesgo
alto (probabilidad de 0.5 o más) y 142 de riesgo medio; 275 no requieren
acción. \codigo{file.exists()} confirma que el archivo quedó en
\archivo{resultados/}. Si el equipo solo puede hacer 150 llamadas, empieza por
los de riesgo alto y sigue en orden de probabilidad (o de pérdida esperada).

¿Qué pasa si alguien le pasa a la función un archivo incompleto? En lugar de un
error críptico de R, recibe un mensaje claro:

\begin{rnoejecutar}
predecir_churn(select(lote_semana, -quejas_6m), modelo_cargado)
#> Error in predecir_churn(select(lote_semana, -quejas_6m), modelo_cargado):
#>   Faltan columnas: quejas_6m
\end{rnoejecutar}

\subsection{Monitoreo y reentrenamiento}

Los modelos se degradan con el tiempo porque el mundo cambia. Se llama
\termino{deriva de datos} (\emph{data drift}) a cuando los datos nuevos ya no
se parecen a los de entrenamiento (por ejemplo, llega una ola de clientes
nuevos por una campaña) y \termino{deriva de concepto} (\emph{concept drift})
a cuando cambia la relación entre variables y respuesta (un competidor lanza
un plan más barato y ahora se van también los clientes antiguos).

Un monitoreo básico compara cada semana las variables del lote con las del
entrenamiento:

\begin{rcode}
comparar_distribucion <- function(entrenamiento, nuevo, variables) {
  map_dfr(variables, function(v) {
    tibble(variable = v,
           media_entrenamiento = mean(entrenamiento[[v]]),
           media_nuevo = mean(nuevo[[v]]))
  }) %>%
    mutate(cambio_pct = round(100 * (media_nuevo - media_entrenamiento) /
                                media_entrenamiento, 1),
           alerta = if_else(abs(cambio_pct) > 20, "REVISAR", "ok"))
}

comparar_distribucion(churn_ent, lote_semana,
                      c("antiguedad_meses", "uso_mensual", "quejas_6m",
                        "pagos_tardios")) %>%
  mutate(across(starts_with("media"), ~ round(.x, 2)))
\end{rcode}

\begin{rsalida}
# A tibble: 4 x 5
  variable         media_entrenamiento media_nuevo cambio_pct alerta
  <chr>                          <dbl>       <dbl>      <dbl> <chr> 
1 antiguedad_meses               36.7        36.8         0.3 ok    
2 uso_mensual                     7.97        8.22        3.2 ok    
3 quejas_6m                       0.78        0.79        2.4 ok    
4 pagos_tardios                   0.41        0.39       -4   ok    
\end{rsalida}

En este lote todas las variables se parecen a las del entrenamiento (cambios
de 4\% o menos), así que no hay alerta. Si un mes la media de quejas saltara
de 0.78 a 1.5 (por ejemplo, por una falla en el servicio), el modelo seguiría
funcionando, pero la tasa de ``Llamar'' se dispararía y convendría revisar
tanto el modelo como la operación. El umbral de 20\% es arbitrario: ajústalo a
tu negocio. Existen medidas más formales, como el índice de estabilidad
poblacional (PSI), pero la idea es la misma.

Buenas prácticas de monitoreo y reentrenamiento:

\begin{itemize}
  \item \textbf{Mide el desempeño real}: tres meses después, compara las
        predicciones con quién realmente se fue. Si el AUC cae de 0.86 a 0.70,
        es hora de reentrenar.
  \item \textbf{Reentrena con calendario} (por ejemplo, cada trimestre) y
        \textbf{por alerta} (cuando el monitoreo detecta cambios).
  \item \textbf{Compara el modelo nuevo contra el actual} en los mismos datos
        de prueba antes de reemplazarlo.
  \item \textbf{Automatiza}: programa el script con el Programador de tareas
        de Windows, \codigo{cron} o RStudio Connect, y envía un reporte
        automático (Módulo~\ref{cap:reportes}).
\end{itemize}

\subsection{Ética, sesgos y explicabilidad}

Un modelo toma decisiones que afectan a personas: a quién se le ofrece un
descuento, a quién se le aprueba un crédito, a quién se contrata. Tienes la
responsabilidad de que sea justo y explicable.

\begin{itemize}
  \item \textbf{Variables sensibles}: género, edad, estado civil, origen o
        código postal (que puede reflejar nivel socioeconómico) pueden llevar
        a discriminar. En nuestro modelo, género y edad no aportaban nada
        (odds ratios de 1.02 y 0.997, importancia mínima): lo correcto es
        \textbf{quitarlas}. En decisiones como crédito o contratación, su uso
        puede ser además ilegal.
  \item \textbf{Variables ``proxy''}: aunque quites el género, otra variable
        muy correlacionada (por ejemplo, la categoría de productos que compra)
        puede reintroducirlo. Revisa las tasas de predicción por grupo.
  \item \textbf{Sesgo histórico}: si en el pasado solo se ofrecían créditos a
        cierto perfil, el modelo aprenderá que ``solo ese perfil paga bien''.
        El modelo repite los sesgos de los datos.
  \item \textbf{Explicabilidad}: prefiere modelos interpretables (logística,
        árboles) cuando la decisión debe justificarse ante un cliente o un
        regulador. Con modelos de caja negra, usa la importancia de variables
        y explica los casos individuales.
  \item \textbf{Privacidad}: usa solo los datos necesarios, respeta los avisos
        de privacidad y la legislación de protección de datos personales.
\end{itemize}

Revisemos si el modelo trata distinto a hombres y mujeres del lote:

\begin{rcode}
puntuados %>%
  group_by(genero) %>%
  summarise(clientes = n(),
            prob_media = round(mean(prob_churn), 3),
            pct_llamar = round(100 * mean(accion == "Llamar"), 1),
            .groups = "drop")
\end{rcode}

\begin{rsalida}
# A tibble: 2 x 4
  genero clientes prob_media pct_llamar
  <chr>     <int>      <dbl>      <dbl>
1 F           258      0.223       40.7
2 M           242      0.238       49.6
\end{rsalida}

La probabilidad media es casi igual (0.223 contra 0.238), pero se llamaría a
49.6\% de los hombres y a 40.7\% de las mujeres. Como el coeficiente de género
es prácticamente nulo, la diferencia viene de las otras variables de este lote
en particular (por azar, los hombres tienen algo más de quejas o menos uso) y
del efecto pequeño pero no nulo que el modelo le asignó a
\codigo{generoM}. Aun así, este tipo de tabla es la que debes revisar y
documentar: si la diferencia fuera grande y se debiera a la variable
sensible, habría que quitarla y volver a evaluar.

\begin{hazlotu}
Reentrena la regresión logística sin \codigo{genero} ni \codigo{edad} y
calcula su AUC en \codigo{churn\_prueba}. ¿Perdiste capacidad predictiva? Si la
respuesta es ``casi nada'', ya tienes un argumento para usar el modelo más
justo y más simple.
\end{hazlotu}

% ============================================================================
\section{Errores comunes}
\label{sec:m7-errores}

\subsection{Fuga de datos: el modelo demasiado bueno}

Agreguemos a los clientes de churn una variable que el área de atención
registra: si el cliente llamó al área de cancelaciones. Como la mayoría de los
que se van llaman para cancelar, parece una variable maravillosa\ldots

\begin{rcode}
set.seed(123)
churn_fuga <- clientes_churn %>%
  mutate(llamo_cancelaciones = if_else(
    churn == "Si", rbinom(n(), 1, 0.85),   # casi todos los que se van
    rbinom(n(), 1, 0.03)))                 # muy pocos de los que se quedan

fuga_ent    <- churn_fuga[filas_ent, ]    # misma división que antes
fuga_prueba <- churn_fuga[-filas_ent, ]

modelo_fuga <- glm(update(formula_churn, . ~ . + llamo_cancelaciones),
                   data = fuga_ent, family = binomial)
prob_fuga <- predict(modelo_fuga, fuga_prueba, type = "response")
round(calcular_auc(prob_fuga, fuga_prueba$churn), 3)
\end{rcode}

\begin{rsalida}
[1] 0.964
\end{rsalida}

El AUC salta de 0.86 a 0.964. En el laboratorio parece un modelo
extraordinario. Pero en la operación no sirve: cuando sabemos que el cliente
llamó a cancelaciones, ya decidió irse. La variable es una consecuencia del
abandono, no una señal anticipada.

\begin{advertencia}
Pregunta obligatoria para cada variable: \textbf{¿conoceré este dato en el
momento en que tenga que predecir?} Si un modelo tiene resultados
``demasiado buenos para ser verdad'' (AUC mayor a 0.95 en un problema de
comportamiento humano), sospecha primero de fuga de datos.
\end{advertencia}

\subsection{Evaluar con los datos de entrenamiento}

Evaluar un modelo con los mismos datos con que se entrenó da resultados
optimistas, sobre todo en modelos flexibles. El random forest es un caso
extremo:

\begin{rcode}
prob_rf_ent <- predict(modelo_rf, newdata = churn_ent,
                       type = "prob")[, "Si"]
c(auc_entrenamiento = calcular_auc(prob_rf_ent, churn_ent$churn),
  auc_prueba = calcular_auc(churn_prueba$prob_rf, churn_prueba$churn)) %>%
  round(3)
\end{rcode}

\begin{rsalida}
auc_entrenamiento        auc_prueba 
            1.000             0.846 
\end{rsalida}

En entrenamiento, el bosque tiene un AUC perfecto (1.000): cada árbol crece
hasta memorizar a sus clientes. En prueba baja a 0.846. Si hubieras reportado
el 1.000, habrías prometido algo imposible. Por eso existen el conjunto de
prueba, la validación cruzada y el error OOB.

\subsection{Clases desbalanceadas y la exactitud engañosa}

Un banco revisa 2{,}000 transacciones con tarjeta, de las cuales 5\% son
fraude. Un ``modelo'' que dice \emph{siempre} ``no es fraude'':

\begin{rcode}
set.seed(123)
fraude <- tibble(
  es_fraude = factor(if_else(runif(2000) < 0.05, "Si", "No"),
                     levels = c("No", "Si"))
)
prediccion_perezosa <- factor(rep("No", 2000), levels = c("No", "Si"))

metricas_clasificacion(prediccion_perezosa, fraude$es_fraude) %>%
  mutate(across(everything(), ~ round(.x, 3)))
\end{rcode}

\begin{rsalida}
# A tibble: 1 x 5
  exactitud precision sensibilidad especificidad    F1
      <dbl>     <dbl>        <dbl>         <dbl> <dbl>
1     0.947       NaN            0             1   NaN
\end{rsalida}

¡94.7\% de exactitud sin detectar un solo fraude! La sensibilidad es 0 y la
precisión y F1 son \codigo{NaN} (\emph{Not a Number}: división entre cero,
porque el modelo nunca predijo ``Si''). Un modelo inútil con una exactitud
altísima: por eso, con clases desbalanceadas, la exactitud es la métrica menos
informativa.

Qué hacer con clases desbalanceadas:

\begin{itemize}
  \item Reporta sensibilidad, precisión, F1 y AUC, no solo exactitud.
  \item Compara siempre con la línea base ``predecir la clase mayoritaria''.
  \item Ajusta el umbral según costos (sección~\ref{sec:m7-umbral}).
  \item Usa divisiones estratificadas (\codigo{createDataPartition()}).
  \item Si hace falta, rebalancea el entrenamiento: \paquete{caret} lo permite
        con \codigo{trainControl(sampling = "down")} o \codigo{"up"}.
\end{itemize}

\subsection{No escalar antes de k-means}

\begin{rcode}
set.seed(123)
kmeans_sin_escalar <- kmeans(rfm %>% select(recencia, frecuencia, monto),
                             centers = 4, nstart = 25)

rfm %>%
  mutate(grupo = kmeans_sin_escalar$cluster) %>%
  group_by(grupo) %>%
  summarise(clientes = n(), recencia = round(mean(recencia)),
            frecuencia = round(mean(frecuencia), 1),
            monto = round(mean(monto)), .groups = "drop") %>%
  arrange(monto)
\end{rcode}

\begin{rsalida}
# A tibble: 4 x 5
  grupo clientes recencia frecuencia monto
  <int>    <int>    <dbl>      <dbl> <dbl>
1     3       56      106        2.9  1946
2     4       83       66        5    4491
3     1       44       39        6.2  7089
4     2       15       37        9.7 11092
\end{rsalida}

Sin escalar, los grupos quedan ordenados casi perfectamente por monto (de
\$1{,}946 a \$11{,}092): el monto, que se mide en miles, domina la distancia.
La recencia apenas participa: ya no aparece un grupo de Dormidos con 195 días
de recencia como en la segmentación escalada, y los clientes que dejaron de
comprar quedan mezclados con los demás según lo que gastaron. Escalar no es un
detalle técnico: cambia los segmentos y, por lo tanto, las estrategias.

\subsection{Otros errores frecuentes}

\begin{table}[H]
\centering
\small
\caption{Otros errores comunes en proyectos de machine learning.}
\label{tab:m7-errores}
\begin{tabularx}{\textwidth}{>{\raggedright\arraybackslash}p{4cm}
                              >{\raggedright\arraybackslash}X}
\toprule
\textbf{Error} & \textbf{Cómo evitarlo} \\
\midrule
Usar un código como número (\codigo{tienda\_id}, código postal) &
  \codigo{lm(ventas \textasciitilde{} tienda\_id)} estima un solo coeficiente, como si la
  tienda 10 fuera ``diez veces'' la tienda 1. Conviértelo con
  \codigo{factor(tienda\_id)}, como hicimos en la sección~\ref{sec:m7-sin-senal}. \\
Niveles de factor distintos entre entrenamiento y datos nuevos &
  Fija los niveles con \codigo{factor(x, levels = ...)} (como en
  \codigo{predecir\_churn()}). Un nivel nuevo (un plan ``Familiar'')
  produce el error \codigo{factor plan has new levels Familiar}. \\
Extrapolar &
  Un modelo entrenado con tiendas de 150 a 1{,}200 m² no sabe nada de una de
  5{,}000 m². Un pronóstico a 5 años con 5 años de historia es una apuesta,
  no una predicción. \\
Separar series de tiempo al azar &
  El conjunto de prueba siempre son los periodos más recientes
  (\codigo{window()}). \\
Olvidar \codigo{set.seed()} &
  Sin semilla, k-means, la división de datos y el random forest dan
  resultados distintos cada vez y nadie puede reproducir tu análisis. \\
Confundir correlación con causalidad &
  Un coeficiente dice ``se asocia con'', no ``provoca''. Para medir el
  efecto de una acción, diseña un experimento (grupo de control). \\
Nombrar clusters por número &
  Asigna nombres a partir del perfil calculado (sección~\ref{sec:m7-segmentacion}). \\
\bottomrule
\end{tabularx}
\end{table}

% ============================================================================
\section{Resumen del módulo}

\begin{resumen}
\begin{itemize}
  \item El machine learning aprende patrones de datos históricos. Supervisado:
        regresión (números) y clasificación (categorías). No supervisado:
        clustering y reglas de asociación. Las series de tiempo pronostican
        con la historia de la propia variable.
  \item Sin señal no hay modelo: primero verifica que exista relación
        (exploración, línea base).
  \item Evalúa siempre en datos que el modelo no vio (prueba, validación
        cruzada); vigila el sobreajuste y la fuga de datos.
  \item Segmentación: RFM + escalar + k-means (codo, silueta, \codigo{nstart})
        + perfil + nombres basados en el perfil + estrategia por segmento.
        Alternativa sencilla: quintiles con \codigo{ntile()}.
  \item Regresión: \codigo{lm()} explica y predice; mide con MAE, RMSE, MAPE y
        R²; compara con una línea base y con random forest.
  \item Clasificación: logística (odds ratios), árbol (reglas), random forest
        (importancia); matriz de confusión, sensibilidad, precisión, AUC; el
        umbral se elige con costos de negocio.
  \item Series de tiempo: \codigo{ts()}, descomposición, prueba con los
        últimos meses, \codigo{naive}/\codigo{snaive} como línea base,
        \codigo{ets()} y \codigo{auto.arima()}; comunica rangos y calcula
        stock de seguridad.
  \item Canasta: soporte, confianza y lift con un self-join; ordena por lift
        con soporte mínimo.
  \item Producción: \codigo{saveRDS()}, función de puntuación con validación,
        exportar resultados, monitorear deriva, reentrenar y cuidar la ética.
\end{itemize}
\end{resumen}

\begin{table}[H]
\centering
\small
\caption{Guía rápida: del problema a la función de R y la métrica.}
\label{tab:m7-resumen}
\begin{tabularx}{\textwidth}{>{\raggedright\arraybackslash}p{3.2cm}
                              >{\raggedright\arraybackslash}p{2.9cm}
                              >{\raggedright\arraybackslash}p{3.9cm}
                              >{\raggedright\arraybackslash}X}
\toprule
\textbf{Problema} & \textbf{Algoritmo} & \textbf{Función de R} & \textbf{Métrica} \\
\midrule
Segmentar clientes & k-means & \codigo{scale()}, \codigo{kmeans()},
  \codigo{fviz\_nbclust()} & Codo (WSS), silueta \\
Segmentar con reglas & Quintiles RFM & \codigo{ntile()}, \codigo{case\_when()} &
  Tamaño y valor de segmentos \\
Predecir un número & Regresión lineal & \codigo{lm()}, \codigo{predict()} &
  MAE, RMSE, MAPE, R² \\
Predecir un número (no lineal) & Random forest & \codigo{randomForest()} &
  MAE, RMSE, OOB \\
Predecir una categoría & Regresión logística & \codigo{glm(family = binomial)} &
  AUC, sensibilidad, precisión, F1 \\
Reglas explicables & Árbol de decisión & \codigo{rpart()} &
  AUC, matriz de confusión \\
Comparar modelos & Validación cruzada & \codigo{trainControl()},
  \codigo{train()}, \codigo{resamples()} & ROC (AUC) promedio \\
Pronosticar demanda & Ingenuo, ETS, ARIMA & \codigo{snaive()}, \codigo{ets()},
  \codigo{auto.arima()}, \codigo{forecast()} & \codigo{accuracy()}: MAE,
  RMSE, MAPE \\
Productos que se compran juntos & Reglas de asociación & self-join con
  \codigo{inner\_join()}; \codigo{arules::apriori()} & Soporte, confianza,
  lift \\
Usar el modelo & Despliegue & \codigo{saveRDS()}, \codigo{readRDS()},
  \codigo{write\_csv()} & Desempeño real y deriva \\
\bottomrule
\end{tabularx}
\end{table}

% ============================================================================
\section{Ejercicios}
\label{sec:m7-ejercicios}

Para los ejercicios con datos simulados, copia el código de simulación de la
sección correspondiente (con la misma semilla) al inicio de tu script. Todos
los ejercicios suponen que tienes cargados los paquetes de este módulo.

\begin{ejercicio}{m7-1}{Supervisado o no supervisado}
\nivel{Básico}
Para cada situación, indica si es un problema de regresión, clasificación,
clustering, reglas de asociación o series de tiempo, y propón un algoritmo:
(a) estimar el ticket promedio de un cliente nuevo; (b) detectar si un correo
de cliente es una queja; (c) agrupar las sucursales por su mezcla de ventas;
(d) predecir la demanda semanal de pan de muerto; (e) decidir qué producto
recomendar a quien compra una cafetera. Después, con
\archivo{ventas\_retail.csv}, ajusta \codigo{lm(total \textasciitilde{} descuento\_pct + dia\_semana)}
y reporta su R² ajustado. ¿Hay señal?
\end{ejercicio}

\begin{ejercicio}{m7-2}{Métricas de regresión a mano}
\nivel{Básico}
Un gerente pronosticó las ventas (miles de pesos) de 6 sucursales:
reales \codigo{c(120, 95, 150, 80, 110, 130)} y pronosticadas
\codigo{c(110, 100, 140, 90, 115, 120)}. Calcula MAE, RMSE, MAPE y R²
escribiendo las fórmulas en R (sin paquetes). ¿Cuál es el error promedio en
pesos y en porcentaje?
\end{ejercicio}

\begin{ejercicio}{m7-3}{Matriz de confusión e interpretación}
\nivel{Básico}
Un modelo de fraude se evaluó con 1{,}000 transacciones: detectó 40 de los 50
fraudes reales y marcó como fraude 60 transacciones legítimas. Construye la
matriz de confusión como un \codigo{matrix()} de R y calcula exactitud,
precisión, sensibilidad, especificidad y F1. Si cada fraude no detectado
cuesta \$5{,}000 y cada revisión innecesaria \$100, ¿cuál es el costo de los
errores?
\end{ejercicio}

\begin{ejercicio}{m7-4}{Segmentación RFM por reglas}
\nivel{Intermedio}
Con \archivo{transacciones.csv} calcula el RFM de cada cliente (fecha de
corte 2024-01-01) y asigna calificaciones de 1 a 5 con \codigo{ntile()}.
Calcula un puntaje total \codigo{r + f + m} y clasifica en ``Oro'' (13 a 15),
``Plata'' (9 a 12) y ``Bronce'' (3 a 8). Reporta, por nivel, número de
clientes, monto promedio y porcentaje de los ingresos totales.
\end{ejercicio}

\begin{ejercicio}{m7-5}{k-means con otras variables}
\nivel{Intermedio}
Con \archivo{transacciones.csv}, construye por cliente: monto total, ticket
promedio (\codigo{mean(total\_transaccion)}), número de tiendas distintas y
número de productos distintos. Escala las cuatro variables, elige $k$ con la
silueta (entre 2 y 6), entrena k-means con \codigo{nstart = 25} y nombra los
grupos a partir de su perfil (sin escribir números de cluster a mano).
\end{ejercicio}

\begin{ejercicio}{m7-6}{Mejorar el modelo de churn con una variable nueva}
\nivel{Intermedio}
Con los datos simulados de churn (misma semilla y misma división de la
sección~\ref{sec:m7-churn}), crea la variable
\codigo{cliente\_nuevo = antiguedad\_meses <= 6}, agrégala a la regresión
logística y compara el AUC en prueba con el del modelo original. Interpreta la
razón de momios de la nueva variable. Quita también \codigo{genero} y
\codigo{edad}: ¿cambia el AUC?
\end{ejercicio}

\begin{ejercicio}{m7-7}{Umbral con otros costos}
\nivel{Intermedio}
Con las probabilidades de la regresión logística en \codigo{churn\_prueba},
encuentra el umbral que minimiza el costo total si el descuento cuesta \$800
(en lugar de \$300), un cliente perdido cuesta \$3{,}000 y la mitad de los
clientes contactados que iban a irse se retiene. ¿Cuántos clientes se
llamarían? Compara con el umbral óptimo del texto (0.15) y explica la
diferencia.
\end{ejercicio}

\begin{ejercicio}{m7-8}{Pronóstico trimestral}
\nivel{Avanzado}
A partir de la serie mensual simulada de la sección~\ref{sec:m7-series},
construye una serie \emph{trimestral} (suma de cada trimestre, frecuencia 4).
Deja como prueba los 4 trimestres de 2025, compara \codigo{snaive()},
\codigo{ets()} y \codigo{auto.arima()} con \codigo{accuracy()} y pronostica
los 4 trimestres de 2026 con el mejor modelo, reentrenado con toda la serie.
\end{ejercicio}

\begin{ejercicio}{m7-9}{Canasta con reglas de tres productos}
\nivel{Avanzado}
Con los tickets simulados de la sección~\ref{sec:m7-canasta}, calcula el
soporte, la confianza y el lift de la regla \{Laptop, Mouse\} $\rightarrow$
\{Funda\} usando \paquete{dplyr}. Compárala con Laptop $\rightarrow$ Funda.
¿Saber que el cliente ya lleva mouse cambia la probabilidad de que lleve funda?
\end{ejercicio}

\begin{ejercicio}{m7-10}{Reto: sistema de puntuación de churn completo}
\nivel{Reto}
Simula (con \codigo{set.seed(2024)}) 3{,}000 clientes con las mismas variables
y reglas de la sección~\ref{sec:m7-churn}. Entrena con \paquete{caret} una
regresión logística y un random forest (validación cruzada de 5 pliegues, AUC
como métrica), elige el mejor según el AUC en un 30\% de prueba, guárdalo con
\codigo{saveRDS()}, escribe una función \codigo{puntuar()} que devuelva los
100 clientes con mayor pérdida esperada (probabilidad $\times$ cuota anual) de
un lote nuevo de 1{,}000 clientes y exporta la lista a
\archivo{resultados/top100\_churn.csv}.
\end{ejercicio}
